\chapter{Planos de conjunto avanzados}\label{cap:planos-conjunto}
% Prefijo de etiquetas: planos-.  (cap:planos ya lo usa el capítulo 6.)

Este capítulo completa los planos del capítulo~\ref{cap:planos} con siete planos de
conjunto: dos vistas isométricas con corte parcial, una vista explosionada con lista de
piezas, la secuencia de despliegue a escala, el recorrido de cargas con sus magnitudes, el
diagrama N$^2$ de interfaces y la envolvente dentro del tubo del cohete. Los planos 9 a 15
siguen la numeración del capítulo~\ref{cap:planos}. Las cotas están en milímetros y las
referencias numeradas coinciden con las de la tabla~\ref{tab:ref-p1} (1 a 19); las nuevas
van de 20 en adelante.

Los cálculos de este capítulo (recorrido de cargas, tiempos de arranque y holguras) están en
\texttt{analysis/planos\_cargas.py}; sus salidas, en \texttt{analysis/data/planos\_*.dat}.
Todos los números de vuelo son \textbf{estimaciones} con el modelo de cantidad de
movimiento del capítulo~\ref{cap:teoria}; ninguno es una medición.

\begin{resultado}[Resumen del capítulo]
\begin{itemize}[leftmargin=*]
  \item El recorrido de cargas es simple y está separado: el tirón del drogue entra por la
        argolla del travesaño y no pasa por los rodamientos; el empuje del rotor
        (\SI{3.8}{N}) sube por los rodamientos hasta la tuerca; el aterrizaje
        (\SI{200}{g} con espuma) es el caso que más carga todas las fijaciones internas
        (\SI{957}{N} en la tapa inferior, \SI{163}{N} en el lastre).
  \item Durante la apertura del drogue la pared queda en \textbf{tracción} debajo del
        travesaño (\SI{34}{N} a \SI{35}{m/s}; \SI{116}{N} con el caso de diseño de 30\,G) y
        el cubo apoya \textbf{hacia abajo} sobre el hombro inferior del eje: la tuerca solo
        trabaja en la etapa 2.
  \item El diagrama N$^2$ identifica 31 interfaces (14 mecánicas, 7 eléctricas y 10 de
        datos). La baliza no tiene interfaz eléctrica con el resto (\Req{E8}).
  \item Las palas plegadas definen hoy el Ø\,\SI{95}{mm}: sus bordes rozarían el tubo del
        cohete. Se propone agregar cuatro patines guía de PETG que definan el Ø\,95 y
        dejar las palas en Ø\,\SI{94}{mm}.
  \item El anexo no da el diámetro interior del tubo. Con la suposición de
        Ø\,\SI{98}{mm} la holgura diametral en el peor caso es de \SI{2.1}{mm}; se
        necesita Ø$_\text{int}\ge\SI{96.9}{mm}$ para conservar \SI{1}{mm}. Hay que
        confirmar el dato con la organización.
\end{itemize}
\end{resultado}

% ---------------------------------------------------------------------
%  Macros de dibujo 3D (prefijo \pl...). Vista: theta = 68, phi = 115.
%  El observador está en el azimut phi - 90 = 25 grados; el arco visible de un
%  cilindro de eje z va de phi - 180 = -65 a phi = 115 grados.
% ---------------------------------------------------------------------
\tdplotsetmaincoords{68}{115}
\def\plA{-65}\def\plB{115}\def\plV{25}
% cilindro macizo: radio, z0, z1, color
\newcommand{\plCil}[4]{%
  \shade[left color=#4!50!white, right color=#4!38!white, middle color=#4!8!white,
         draw=#4!75!black, line width=0.35pt]
    plot[domain=\plA:\plB, samples=40, variable=\t] ({#1*cos(\t)},{#1*sin(\t)},{#2})
    -- plot[domain=\plB:\plA, samples=40, variable=\t] ({#1*cos(\t)},{#1*sin(\t)},{#3})
    -- cycle;
  \filldraw[fill=#4!20!white, draw=#4!75!black, line width=0.35pt]
    plot[domain=0:360, samples=73, variable=\t] ({#1*cos(\t)},{#1*sin(\t)},{#3}) -- cycle;}
% tubo o anillo: radio ext., radio int., z0, z1, color
\newcommand{\plTubo}[5]{%
  \plCil{#1}{#3}{#4}{#5}%
  \filldraw[fill=#5!55!black!45, draw=#5!75!black, line width=0.3pt]
    plot[domain=0:360, samples=73, variable=\t] ({#2*cos(\t)},{#2*sin(\t)},{#4}) -- cycle;}
% faja cilíndrica (pala plegada, trozo de pared): radio, a1, a2, z0, z1, color
\newcommand{\plFaja}[6]{%
  \shade[left color=#6!45!white, right color=#6!30!white, middle color=#6!10!white,
         draw=#6!75!black, line width=0.35pt]
    plot[domain=#2:#3, samples=25, variable=\t] ({#1*cos(\t)},{#1*sin(\t)},{#4})
    -- plot[domain=#3:#2, samples=25, variable=\t] ({#1*cos(\t)},{#1*sin(\t)},{#5})
    -- cycle;}
% caja radial: azimut, r1, r2, semiancho, z0, z1, color
\newcommand{\plCaja}[7]{%
  \pgfmathsetmacro{\plca}{cos(#1)}\pgfmathsetmacro{\plsa}{sin(#1)}%
  \pgfmathtruncatemacro{\plvo}{cos(#1-\plV)>0}%
  \pgfmathtruncatemacro{\plvs}{sin(\plV-#1)>0}%
  \def\plw{#4}\ifnum\plvs=0 \def\plw{-#4}\fi
  \filldraw[fill=#7!32!white, draw=#7!75!black, line width=0.3pt]
    ({#2*\plca-\plw*\plsa},{#2*\plsa+\plw*\plca},{#5}) --
    ({#3*\plca-\plw*\plsa},{#3*\plsa+\plw*\plca},{#5}) --
    ({#3*\plca-\plw*\plsa},{#3*\plsa+\plw*\plca},{#6}) --
    ({#2*\plca-\plw*\plsa},{#2*\plsa+\plw*\plca},{#6}) -- cycle;
  \ifnum\plvo=1
  \filldraw[fill=#7!45!white, draw=#7!75!black, line width=0.3pt]
    ({#3*\plca+#4*\plsa},{#3*\plsa-#4*\plca},{#5}) --
    ({#3*\plca-#4*\plsa},{#3*\plsa+#4*\plca},{#5}) --
    ({#3*\plca-#4*\plsa},{#3*\plsa+#4*\plca},{#6}) --
    ({#3*\plca+#4*\plsa},{#3*\plsa-#4*\plca},{#6}) -- cycle;
  \else
  \filldraw[fill=#7!45!white, draw=#7!75!black, line width=0.3pt]
    ({#2*\plca+#4*\plsa},{#2*\plsa-#4*\plca},{#5}) --
    ({#2*\plca-#4*\plsa},{#2*\plsa+#4*\plca},{#5}) --
    ({#2*\plca-#4*\plsa},{#2*\plsa+#4*\plca},{#6}) --
    ({#2*\plca+#4*\plsa},{#2*\plsa-#4*\plca},{#6}) -- cycle;
  \fi
  \filldraw[fill=#7!15!white, draw=#7!75!black, line width=0.3pt]
    ({#2*\plca+#4*\plsa},{#2*\plsa-#4*\plca},{#6}) --
    ({#3*\plca+#4*\plsa},{#3*\plsa-#4*\plca},{#6}) --
    ({#3*\plca-#4*\plsa},{#3*\plsa+#4*\plca},{#6}) --
    ({#2*\plca-#4*\plsa},{#2*\plsa+#4*\plca},{#6}) -- cycle;}
% cara de corte radial (rayada): azimut, r1, r2, z0, z1, color
\newcommand{\plCorte}[6]{%
  \filldraw[fill=#6!25!white, draw=#6!80!black, line width=0.3pt]
    ({#2*cos(#1)},{#2*sin(#1)},{#4}) -- ({#3*cos(#1)},{#3*sin(#1)},{#4}) --
    ({#3*cos(#1)},{#3*sin(#1)},{#5}) -- ({#2*cos(#1)},{#2*sin(#1)},{#5}) -- cycle;
  \fill[pattern=north east lines, pattern color=#6!80]
    ({#2*cos(#1)},{#2*sin(#1)},{#4}) -- ({#3*cos(#1)},{#3*sin(#1)},{#4}) --
    ({#3*cos(#1)},{#3*sin(#1)},{#5}) -- ({#2*cos(#1)},{#2*sin(#1)},{#5}) -- cycle;}
% llamada: punto, posición del globo (coordenadas de lienzo), número
\newcommand{\plG}[3]{\draw[ref] (#1) -- (#2); \node[marca] at (#2) {#3};}

%=====================================================================
\section{Plano 9: vistas isométricas del CanSat plegado}\label{sec:planos-iso}

La vista isométrica muestra en un solo dibujo cómo se apilan las piezas y dónde quedan las
interfaces que el corte longitudinal del plano~1 deja ocultas. La
figura~\ref{fig:planos-iso-plegado} tiene dos vistas a escala 1:2: el exterior tal como
entra al cohete y un corte parcial que quita un sector de \ang{90} del cuerpo, de la tapa
superior y de las dos palas delanteras.

\begin{figure}[p]
\centering
\begin{subfigure}[b]{0.53\textwidth}\centering
\begin{tikzpicture}[scale=0.05]
\begin{scope}[tdplot_main_coords]
  % cuerpo
  \plCil{44}{0}{205}{cfijo}
  % orificio de interruptores en el hueco delantero (azimut 25)
  \draw[cfijo!70!black, fill=white, line width=0.3pt]
    plot[domain=0:360, samples=30, variable=\s]
      ({44*cos(25+5.2*cos(\s))},{44*sin(25+5.2*cos(\s))},{120+4*sin(\s)});
  \draw[cfijo!70!black, fill=white, line width=0.3pt]
    plot[domain=0:360, samples=30, variable=\s]
      ({44*cos(25+5.2*cos(\s))},{44*sin(25+5.2*cos(\s))},{100+4*sin(\s)});
  \fill[ccuerda!50!black] ({44*cos(25)},{44*sin(25)},{140}) circle (1.6);
  % palas delanteras (azimut -20 y 70), plegadas
  \plFaja{45.6}{-48.25}{8.25}{57}{211}{crot}
  \plFaja{45.6}{41.75}{98.25}{57}{211}{crot}
  % pestañas de punta
  \foreach \a in {-20,70}{\fill[black] ({46*cos(\a)},{46*sin(\a)},{58.5}) circle (1.3);}
  % brazos y portapalas traseros
  \plCaja{160}{22}{46.5}{4}{209}{217}{crot}
  \plCaja{-110}{22}{46.5}{4}{209}{217}{crot}
  % cubo
  \plCil{22}{205.5}{222}{crot}
  \fill[black] ({18*cos(40)},{18*sin(40)},{222}) circle (1.2);
  % portapalas y brazos delanteros
  \plCaja{-20}{44.3}{46.8}{6}{192}{211}{crot}
  \plCaja{70}{44.3}{46.8}{6}{192}{211}{crot}
  \plCaja{-20}{22}{46.5}{4}{209}{217}{crot}
  \plCaja{70}{22}{46.5}{4}{209}{217}{crot}
  % tuerca, copa y cuerda
  \plCil{7}{222}{225}{cgris}
  \plTubo{22}{20.5}{225}{250}{cfijo}
  \fill[ccuerda!35] plot[domain=0:360, samples=50, variable=\t]
      ({20.5*cos(\t)},{20.5*sin(\t)},{247}) -- cycle;
  \draw[cuerda] (0,0,247) -- (0,0,272);
  % puntos de llamada
  \coordinate (a1) at (0,0,268);
  \coordinate (a2) at ({22*cos(80)},{22*sin(80)},{240});
  \coordinate (a4) at ({22*cos(80)},{22*sin(80)},{214});
  \coordinate (a7) at ({46.8*cos(76)},{46.8*sin(76)},{200});
  \coordinate (a8) at ({45.6*cos(85)},{45.6*sin(85)},{140});
  \coordinate (a9) at ({44*cos(40)},{44*sin(40)},{30});
  \coordinate (a14) at ({46*cos(70)},{46*sin(70)},{58.5});
  \coordinate (a19) at ({44*cos(25)},{44*sin(25)},{120});
  \coordinate (a28) at ({44*cos(25)},{44*sin(25)},{140});
  \coordinate (a6) at ({45.2*cos(-20)},{45.2*sin(-20)},{213});
  \coordinate (b0) at (0,0,0);
  \coordinate (b250) at (0,0,250);
  \coordinate (bR) at ({47.5*cos(115)},{47.5*sin(115)},{0});
  \coordinate (bL) at ({47.5*cos(-65)},{47.5*sin(-65)},{0});
\end{scope}
  \path (-76,0) coordinate (colL) (80,0) coordinate (colR);
  \plG{a1}{a1 -| colR}{1}
  \plG{a2}{$(a2 -| colR)+(0,4)$}{2}
  \plG{a4}{$(a4 -| colR)+(0,-6)$}{4}
  \plG{a6}{$(a6 -| colL)+(0,4)$}{6}
  \plG{a7}{$(a7 -| colR)+(0,-14)$}{7}
  \plG{a8}{a8 -| colR}{8}
  \plG{a9}{a9 -| colR}{9}
  \plG{a14}{a14 -| colR}{14}
  \plG{a19}{$(a19 -| colL)+(0,-6)$}{19}
  \plG{a28}{$(a28 -| colL)+(0,6)$}{28}
  % cota de largo
  \draw[cgris, line width=0.3pt] (b0) -- ++(-88,0) (b250) -- ++(-88,0);
  \draw[cota] ($(b0)+(-84,0)$) -- ($(b250 -| b0)+(-84,0)$)
     node[midway, fill=white, font=\scriptsize, rotate=90] {250 (\Req{S3})};
\end{tikzpicture}
\caption{Exterior, listo para integrar.}
\end{subfigure}\hfill
\begin{subfigure}[b]{0.45\textwidth}\centering
\begin{tikzpicture}[scale=0.05]
\begin{scope}[tdplot_main_coords]
  % pared interior trasera visible por el corte
  \fill[cfijo!30!white]
    plot[domain=115:295, samples=40, variable=\t] ({42*cos(\t)},{42*sin(\t)},{3})
    -- plot[domain=295:115, samples=40, variable=\t] ({42*cos(\t)},{42*sin(\t)},{199}) -- cycle;
  % piso (cara superior de la tapa inferior)
  \fill[cfijo!12!white, draw=cfijo!60, line width=0.3pt]
    plot[domain=0:360, samples=60, variable=\t] ({42*cos(\t)},{42*sin(\t)},{3}) -- cycle;
  % baliza y cámara nadir sobre la tapa inferior
  \begin{scope}[shift={(-12,-24,0)}]\plCil{6}{3}{16}{cfijo}\end{scope}
  \begin{scope}[shift={(20,10,0)}]\plCil{8}{3}{14}{cgris}\end{scope}
  % lastre
  \plCil{36}{20}{30}{cgris}
  % servo y anillo de traba
  \plTubo{35}{31}{52}{58}{cfreno}
  \plCaja{25}{-13}{13}{8}{40}{64}{cfreno}
  % batería 18650
  \plCil{9}{72}{137}{cfijo}
  % separadores y placas
  \foreach \a in {205,295}{\begin{scope}[shift={({34*cos(\a)},{34*sin(\a)},0)}]\plCil{1.5}{140}{172}{cgris}\end{scope}}
  \foreach \zp in {146,158,170}{\plCil{38}{\zp-1.6}{\zp}{ccuerda!70!black}}
  \foreach \a in {25,115}{\begin{scope}[shift={({34*cos(\a)},{34*sin(\a)},0)}]\plCil{1.5}{140}{172}{cgris}\end{scope}}
  % travesaño y argolla
  \plCaja{-20}{-20}{20}{4}{176}{180}{cfijo}
  % eje fijo y cuerda dentro del cuerpo
  \plCil{4}{180}{205}{cgris}
  % cámara cenit (fuera del eje)
  \begin{scope}[shift={({28*cos(25)},{28*sin(25)},0)}]\plCil{5}{188}{199}{cfijo}\end{scope}
  % cascarón exterior fuera del sector cortado
  \plFaja{44}{-65}{-20}{0}{205}{cfijo}
  \plFaja{44}{70}{115}{0}{205}{cfijo}
  % tapa superior con sector quitado
  \filldraw[fill=cfijo!20!white, draw=cfijo!75!black, line width=0.35pt]
    (0,0,205) -- plot[domain=70:340, samples=50, variable=\t] ({44*cos(\t)},{44*sin(\t)},{205}) -- cycle;
  % caras de corte
  \plCorte{-20}{42}{44}{0}{205}{cfijo}
  \plCorte{70}{42}{44}{0}{205}{cfijo}
  \plCorte{-20}{0}{42}{0}{3}{cfijo}
  \plCorte{70}{0}{42}{0}{3}{cfijo}
  \plCorte{-20}{4}{42}{199}{205}{cfijo}
  \plCorte{70}{4}{42}{199}{205}{cfijo}
  % palas: solo las mitades fuera del corte
  \plFaja{45.6}{-48.25}{-20}{57}{211}{crot}
  \plFaja{45.6}{70}{98.25}{57}{211}{crot}
  % rotor
  \plCaja{160}{22}{46.5}{4}{209}{217}{crot}
  \plCaja{-110}{22}{46.5}{4}{209}{217}{crot}
  \plCil{22}{205.5}{222}{crot}
  \plCaja{-20}{22}{46.5}{4}{209}{217}{crot}
  \plCaja{70}{22}{46.5}{4}{209}{217}{crot}
  \plCil{7}{222}{225}{cgris}
  \plTubo{22}{20.5}{225}{250}{cfijo}
  \fill[ccuerda!35] plot[domain=0:360, samples=50, variable=\t]
      ({20.5*cos(\t)},{20.5*sin(\t)},{247}) -- cycle;
  \draw[cuerda] (0,0,247) -- (0,0,272);
  % puntos de llamada
  \coordinate (c12) at (0,0,192);
  \coordinate (c13) at ({9*cos(-20)},{9*sin(-20)},{105});
  \coordinate (c15) at ({35*cos(-10)},{35*sin(-10)},{55});
  \coordinate (c16) at ({36*cos(30)},{36*sin(30)},{25});
  \coordinate (c17) at (-12,-24,12);
  \coordinate (c18) at (20,10,10);
  \coordinate (c24) at ({18*cos(-20)},{18*sin(-20)},{178});
  \coordinate (c25) at ({38*cos(0)},{38*sin(0)},{158});
  \coordinate (c23) at ({30*cos(-20)},{30*sin(-20)},{202});
  \coordinate (c9) at ({43*cos(-20)},{43*sin(-20)},{120});
  \coordinate (c11) at ({28*cos(25)},{28*sin(25)},{196});
\end{scope}
  \path (-64,0) coordinate (colL) (80,0) coordinate (colR);
  \plG{c12}{$(c12 -| colR)+(0,22)$}{12}
  \plG{c23}{$(c23 -| colR)+(0,4)$}{23}
  \plG{c24}{$(c24 -| colR)+(0,-4)$}{24}
  \plG{c25}{c25 -| colR}{25}
  \plG{c13}{c13 -| colR}{13}
  \plG{c15}{c15 -| colR}{15}
  \plG{c16}{$(c16 -| colR)+(0,-6)$}{16}
  \plG{c17}{$(c17 -| colR)+(0,10)$}{17}
  \plG{c18}{$(c18 -| colR)+(0,-14)$}{18}
  \plG{c9}{c9 -| colL}{9}
  \plG{c11}{$(c11 -| colL)+(0,10)$}{11}
  \rotulo{$(current bounding box.south east)+(0,-30)$}{Plano 9 -- CanSat plegado}{1:2}{9}
\end{tikzpicture}
\caption{Corte parcial: sector de \ang{90} retirado.}
\end{subfigure}
\caption{Plano 9: vistas isométricas del CanSat plegado, configuración A (escala 1:2;
vista desde \ang{22} sobre el horizonte). Azul: piezas fijas; naranja: rotor; rojo: traba;
verde: drogue, cuerda y placas de circuito. La cámara nadir (18) y la baliza (17) apoyan en
la tapa inferior; el lastre (16) queda justo encima. Referencias en la
tabla~\ref{tab:planos-piezas}.}
\label{fig:planos-iso-plegado}
\end{figure}

En la vista exterior se ve que los cuatro huecos de \ang{33.5} entre palas (azimut del
observador) son el único acceso al cuerpo con el CanSat integrado: allí van los orificios
de los interruptores (19, \Req{E4}, \Req{E5}, \Req{E9}) y el indicador de encendido (28,
\Req{E6}). El corte muestra el orden de la pila interna de abajo hacia arriba: baliza y
cámara nadir, lastre, anillo de traba con su servo, batería vertical, tres placas sobre
separadores, travesaño de anclaje y eje fijo. Las palas plegadas son exteriores al
cuerpo y ocupan la franja de \SIrange{57}{211}{mm}; ninguna pieza interna las cruza.

%=====================================================================
\section{Plano 10: vista isométrica del CanSat desplegado}\label{sec:planos-iso-desp}

Desplegado, el rotor ocupa un disco de Ø\,\SI{400}{mm}, cuatro veces el diámetro del
cuerpo; el drogue (Ø\,\SI{250}{mm}) queda \SI{1.5}{m} más arriba, fuera de la estela
del rotor según la discusión del capítulo~\ref{cap:teoria}. La
figura~\ref{fig:planos-iso-desp} lo dibuja a escala 1:3,5, con la cuerda cortada.

\newcommand{\plPala}[2]{% azimut, color
  \pgfmathsetmacro{\pca}{cos(#1)}\pgfmathsetmacro{\psa}{sin(#1)}%
  \def\plZ##1##2{213+(##1-45.2)*0.0769-(##2)*0.0524}%
  \filldraw[fill=#2!25!white, draw=#2!80!black, line width=0.35pt]
    ({46*\pca+22.5*\psa},{46*\psa-22.5*\pca},{\plZ{46}{-22.5}}) --
    ({200*\pca+22.5*\psa},{200*\psa-22.5*\pca},{\plZ{200}{-22.5}}) --
    ({200*\pca-22.5*\psa},{200*\psa+22.5*\pca},{\plZ{200}{22.5}}) --
    ({46*\pca-22.5*\psa},{46*\psa+22.5*\pca},{\plZ{46}{22.5}}) -- cycle;
  \fill[#2!45!white]
    ({46*\pca+22.5*\psa},{46*\psa-22.5*\pca},{\plZ{46}{-22.5}}) --
    ({66*\pca+22.5*\psa},{66*\psa-22.5*\pca},{\plZ{66}{-22.5}}) --
    ({66*\pca-22.5*\psa},{66*\psa+22.5*\pca},{\plZ{66}{22.5}}) --
    ({46*\pca-22.5*\psa},{46*\psa+22.5*\pca},{\plZ{46}{22.5}}) -- cycle;
  \draw[#2, line width=1.3pt]
    ({46*\pca-22.5*\psa},{46*\psa+22.5*\pca},{\plZ{46}{22.5}}) --
    ({200*\pca-22.5*\psa},{200*\psa+22.5*\pca},{\plZ{200}{22.5}});}

\begin{figure}[p]
\centering
\begin{tikzpicture}[scale=0.0286]
\begin{scope}[tdplot_main_coords]
  % trayectoria de puntas
  \draw[crot, densely dashed, line width=0.4pt]
    plot[domain=0:360, samples=90, variable=\t] ({200*cos(\t)},{200*sin(\t)},{225});
  % cuerpo
  \plCil{44}{0}{205}{cfijo}
  % palas traseras, cubo, palas delanteras
  \plPala{160}{crot}\plPala{-110}{crot}
  \plCaja{160}{22}{46.5}{4}{209}{217}{crot}
  \plCaja{-110}{22}{46.5}{4}{209}{217}{crot}
  \plCil{22}{205.5}{222}{crot}
  \plCaja{-20}{22}{46.5}{4}{209}{217}{crot}
  \plCaja{70}{22}{46.5}{4}{209}{217}{crot}
  \plPala{-20}{crot}\plPala{70}{crot}
  \plCil{7}{222}{225}{cgris}
  \plTubo{22}{20.5}{225}{250}{cfijo}
  % cuerda (cortada) y drogue
  \draw[cuerda] (0,0,248) -- (0,0,300);
  \draw[cuerda] (0,0,330) -- (0,0,372);
  \coordinate (k1) at (0,0,300); \coordinate (k2) at (0,0,330);
  \fill[ccuerda] (0,0,372) circle (3);
  \foreach \c in {15,45,...,345}{
    \draw[ccuerda!80, line width=0.25pt] (0,0,372) -- ({125*cos(\c-15)},{125*sin(\c-15)},{430});}
  \foreach \c in {195,225,165,255,135,285,105,315,75,345,45,15}{
    \pgfmathtruncatemacro{\k}{mod((\c-15)/30,2)}
    \ifnum\k=0 \def\gc{ccuerda!22!white}\else\def\gc{ccuerda!50!white}\fi
    \filldraw[fill=\gc, draw=ccuerda!70!black, line width=0.3pt]
      plot[domain=0:90, samples=12, variable=\s]
        ({125*cos(\s)*cos(\c-15)},{125*cos(\s)*sin(\c-15)},{430+48*sin(\s)})
      -- plot[domain=90:0, samples=12, variable=\s]
        ({125*cos(\s)*cos(\c+15)},{125*cos(\s)*sin(\c+15)},{430+48*sin(\s)}) -- cycle;}
  % sentido de giro
  \draw[crot!80!black, -{Latex[length=2.5mm]}, line width=0.9pt]
    plot[domain=-35:35, samples=30, variable=\t] ({218*cos(\t+115)},{218*sin(\t+115)},{232});
  % puntos de referencia
  \coordinate (d0) at ({200*cos(-65)},{200*sin(-65)},{225});
  \coordinate (d1) at ({200*cos(115)},{200*sin(115)},{225});
  \coordinate (h0) at (0,0,0);
  \coordinate (h1) at (0,0,250);
  \coordinate (pala) at ({140*cos(-20)},{140*sin(-20)},{220});
  \coordinate (ba) at ({170*cos(70)-22.5*sin(70)},{170*sin(70)+22.5*cos(70)},{222});
  \coordinate (cubo) at ({22*cos(25)},{22*sin(25)},{214});
  \coordinate (dro) at ({125*cos(-65)},{125*sin(-65)},{432});
  \coordinate (cue) at (0,0,350);
  \coordinate (cpo) at ({44*cos(60)},{44*sin(60)},{100});
  \coordinate (giro) at ({218*cos(80)},{218*sin(80)},{232});
\end{scope}
  % rotura de la cuerda
  \draw[black, line width=0.4pt] ($(k1)+(-8,-3)$) -- ($(k1)+(8,3)$) ($(k2)+(-8,-3)$) -- ($(k2)+(8,3)$);
  \node[font=\scriptsize, anchor=east, align=right] at ($(k1)!0.5!(k2)+(-10,0)$)
       {cuerda Dyneema, \SI{1.5}{m}\\ (cortada en el dibujo)};
  % cotas
  \draw[cota] ($(d0 |- h0)+(0,-40)$) -- ($(d1 |- h0)+(0,-40)$)
     node[pos=0.22, fill=white, font=\scriptsize] {Ø\,400 ($R=200$)};
  \draw[cgris, line width=0.3pt] (d0) -- ($(d0 |- h0)+(0,-46)$) (d1) -- ($(d1 |- h0)+(0,-46)$);
  % notas
  \node[font=\small, crot!80!black, anchor=south west] at ($(giro)+(8,8)$) {$\Omega$};
  \draw[ref] (dro) -- ++(-40,25) node[anchor=east, font=\scriptsize, align=right]
     {drogue Ø\,250, $C_D\approx0{,}8$\\ (sigue atado en la etapa 2)};
  \draw[ref] (pala) -- ++(-60,-120) node[anchor=north east, font=\scriptsize, align=right]
     {pala 156 × 45, conicidad $\beta_0\approx\ang{4.4}$\\ paso $\theta\approx\ang{-3}$};
  \draw[ref] (ba) -- ++(40,-100) node[anchor=north west, font=\scriptsize, align=left]
     {borde de ataque\\ (línea gruesa)};
  \draw[ref] (cpo) -- ++(70,-60) node[anchor=west, font=\scriptsize, align=left]
     {cuerpo Ø\,88\\ $z_\text{cg}\approx\SI{107}{mm}$};
  \draw[flecha, cgris] ($(h0)+(-235,300)$) -- ++(0,-110)
     node[midway, left, font=\scriptsize, align=right] {$V_2\approx\SI{6.4}{m/s}$\\ (estimado)};
  \rotulo{$(current bounding box.south east)+(-45,-60)$}{Plano 10 -- CanSat desplegado}{1:3,5}{10}
\end{tikzpicture}
\caption{Plano 10: CanSat desplegado en la etapa 2, vista isométrica a escala 1:3,5 (cuerda
del drogue cortada). El rotor gira en sentido antihorario visto desde arriba ($\Omega\approx\SI{1150}{rpm}$,
estimado). Palas dibujadas planas con la conicidad y el paso de la base; la
curvatura del 7\,\% no se ve a esta escala. El capítulo~\ref{cap:estructura} estima una
conicidad real de $\approx\ang{6.1}$ con la distribución de masa real; la diferencia no
cambia la envolvente.}
\label{fig:planos-iso-desp}
\end{figure}

%=====================================================================
\section{Plano 11: vista explosionada y lista de piezas}\label{sec:planos-explo}

La vista explosionada ordena el montaje en tres grupos que se arman por separado y se unen
al final: la cabeza del rotor (grupo A), el cuerpo (grupo B) y el módulo interno con la
electrónica y el lastre (grupo C). La figura~\ref{fig:planos-explo} los muestra a escala
1:2,5 sobre sus ejes de montaje, y la tabla~\ref{tab:planos-piezas} lista las 30 piezas.

\begin{figure}[p]
\centering
\begin{tikzpicture}[scale=0.038]
% ------------------------------------------------------------ grupo A: cabeza del rotor
\begin{scope}[tdplot_main_coords]
  \draw[eje] (0,0,95) -- (0,0,490);
  % travesaño y argolla
  \plCaja{-20}{-20}{20}{4}{110}{114}{cfijo}
  \fill[ccuerda] (0,0,108) circle (2.4);
  % tapa superior con sensor hall y cámara cenit
  \plTubo{44}{4}{140}{146}{cfijo}
  \begin{scope}[shift={({15.5*cos(-20)},{15.5*sin(-20)},0)}]\plCil{2}{146}{148}{cfreno}\end{scope}
  \begin{scope}[shift={({28*cos(25)},{28*sin(25)},0)}]\plCil{5}{146}{157}{cfijo}\end{scope}
  % eje fijo
  \plTubo{4}{2.5}{168}{235}{cgris}
  % cubo y brazos
  \plCaja{160}{22}{46.5}{4}{259}{267}{crot}
  \plCaja{-110}{22}{46.5}{4}{259}{267}{crot}
  \plTubo{22}{8}{255}{271.5}{crot}
  \fill[black] ({16*cos(40)},{16*sin(40)},{271.5}) circle (1.3);
  \plCaja{-20}{22}{46.5}{4}{259}{267}{crot}
  \plCaja{70}{22}{46.5}{4}{259}{267}{crot}
  % pala con portapala, separada del brazo delantero derecho
  \pgfmathsetmacro{\eca}{cos(70)}\pgfmathsetmacro{\esa}{sin(70)}
  \filldraw[fill=frot, draw=crot!80!black, line width=0.35pt]
     ({92*\eca+22.5*\esa},{92*\esa-22.5*\eca},{90}) --
     ({92*\eca-22.5*\esa},{92*\esa+22.5*\eca},{90}) --
     ({92*\eca-22.5*\esa},{92*\esa+22.5*\eca},{246}) --
     ({92*\eca+22.5*\esa},{92*\esa-22.5*\eca},{246}) -- cycle;
  \draw[crot, line width=1.2pt]
     ({92*\eca-22.5*\esa},{92*\esa+22.5*\eca},{90}) --
     ({92*\eca-22.5*\esa},{92*\esa+22.5*\eca},{246});
  \plCaja{70}{90.7}{93.3}{6}{227}{246}{crot}
  \fill[black] ({92*\eca},{92*\esa},{91}) circle (1.4);
  % rodamientos y separador
  \plTubo{8}{4}{300}{305}{cgris}
  \plTubo{6}{4}{318}{321}{cgris}
  \plTubo{8}{4}{336}{341}{cgris}
  % tuerca
  \plTubo{7}{4}{360}{363}{cgris}
  % copa
  \plTubo{22}{20.5}{385}{410}{cfijo}
  % drogue plegado y cuerda
  \plCil{17}{430}{452}{ccuerda}
  \draw[cuerda] (0,0,452) -- (0,0,480);
  \coordinate (e1) at (0,0,472); \coordinate (e20) at ({17*cos(-20)},{17*sin(-20)},{441});
  \coordinate (e2) at ({22*cos(-30)},{22*sin(-30)},{398});
  \coordinate (e3) at ({7*cos(-40)},{7*sin(-40)},{361.5});
  \coordinate (e5a) at ({8*cos(-40)},{8*sin(-40)},{338.5});
  \coordinate (e22) at ({6*cos(-40)},{6*sin(-40)},{319.5});
  \coordinate (e5b) at ({8*cos(-40)},{8*sin(-40)},{302.5});
  \coordinate (e4) at ({22*cos(-40)},{22*sin(-40)},{263});
  \coordinate (e10) at ({16*cos(40)},{16*sin(40)},{271.5});
  \coordinate (e7) at ({93.3*cos(70)},{93.3*sin(70)},{238});
  \coordinate (e8) at ({92*cos(70)},{92*sin(70)},{160});
  \coordinate (e14) at ({92*cos(70)},{92*sin(70)},{91});
  \coordinate (e6) at ({46.5*cos(70)},{46.5*sin(70)},{263});
  \coordinate (e12) at ({4*cos(-30)},{4*sin(-30)},{200});
  \coordinate (e23) at ({44*cos(-40)},{44*sin(-40)},{143});
  \coordinate (e11) at ({28*cos(25)},{28*sin(25)},{153});
  \coordinate (e24) at ({20*cos(-20)},{20*sin(-20)},{112});
\end{scope}
% ------------------------------------------------------------ grupo B: cuerpo
\begin{scope}[shift={(165,0)}]
\begin{scope}[tdplot_main_coords]
  \draw[eje] (0,0,95) -- (0,0,400);
  \plTubo{44}{42}{150}{355}{cfijo}
  % patines guía (propuestos), en los huecos entre palas
  \foreach \a in {25,115}{\plCaja{\a}{44}{47.5}{2}{150}{355}{ccont}}
  % orificios
  \foreach \zz in {250,270}{
  \draw[cfijo!70!black, fill=white, line width=0.3pt]
    plot[domain=0:360, samples=30, variable=\s]
      ({44*cos(40+5.2*cos(\s))},{44*sin(40+5.2*cos(\s))},{\zz+4*sin(\s)});}
  \fill[cfreno] ({44*cos(40)},{44*sin(40)},{290}) circle (1.6);
  % ranuras de las pestañas a z = 60
  \foreach \a in {-20,70}{\fill[black] ({44*cos(\a)},{44*sin(\a)},{210}) circle (1.3);}
  \coordinate (e9) at ({44*cos(-50)},{44*sin(-50)},{300});
  \coordinate (e27) at ({47.5*cos(25)},{47.5*sin(25)},{200});
  \coordinate (e19) at ({44*cos(40)},{44*sin(40)},{250});
  \coordinate (e28) at ({44*cos(40)},{44*sin(40)},{290});
  \coordinate (r14) at ({44*cos(70)},{44*sin(70)},{210});
\end{scope}
\end{scope}
% ------------------------------------------------------------ grupo C: módulo interno
\begin{scope}[shift={(300,0)}]
\begin{scope}[tdplot_main_coords]
  \draw[eje] (0,0,30) -- (0,0,380);
  % tapa inferior
  \plCil{44}{40}{43}{cfijo}
  % baliza y cámara nadir
  \begin{scope}[shift={(-12,-24,0)}]\plCil{6}{63}{76}{cfijo}\end{scope}
  \begin{scope}[shift={(20,10,0)}]\plCil{8}{63}{74}{cgris}\end{scope}
  % lastre
  \plCil{36}{100}{110}{cgris}
  % anillo de traba y servo
  \plCaja{25}{-13}{13}{8}{135}{159}{cfreno}
  \plTubo{35}{31}{170}{176}{cfreno}
  % batería
  \plCil{9}{200}{265}{cfijo}
  % separadores y placas
  \foreach \a in {205,295}{\begin{scope}[shift={({34*cos(\a)},{34*sin(\a)},0)}]\plCil{1.5}{290}{350}{cgris}\end{scope}}
  \foreach \zp in {300,320,340}{\plCil{38}{\zp-1.6}{\zp}{ccuerda!70!black}}
  \foreach \a in {25,115}{\begin{scope}[shift={({34*cos(\a)},{34*sin(\a)},0)}]\plCil{1.5}{290}{350}{cgris}\end{scope}}
  \coordinate (e26) at ({44*cos(-40)},{44*sin(-40)},{41.5});
  \coordinate (e17) at (-12,-24,70);
  \coordinate (e18) at (20,10,70);
  \coordinate (e16) at ({36*cos(-40)},{36*sin(-40)},{105});
  \coordinate (e15) at ({35*cos(-40)},{35*sin(-40)},{173});
  \coordinate (e15b) at (0,0,147);
  \coordinate (e13) at ({9*cos(-40)},{9*sin(-40)},{232});
  \coordinate (e25) at ({38*cos(-40)},{38*sin(-40)},{320});
  \coordinate (e30) at ({50*cos(25)},{50*sin(25)},{400});
\end{scope}
\end{scope}
% ------------------------------------------------------------ flechas de montaje y globos
  \path (-75,0) coordinate (cA) (96,0) coordinate (cA2) (250,0) coordinate (cB2) (378,0) coordinate (cC);
  \plG{e1}{e1 -| cA}{1}
  \plG{e20}{e20 -| cA}{20}
  \plG{e2}{e2 -| cA}{2}
  \plG{e3}{$(e3 -| cA)+(0,4)$}{3}
  \plG{e5a}{$(e5a -| cA)+(0,-2)$}{5}
  \plG{e22}{e22 -| cA}{22}
  \plG{e5b}{$(e5b -| cA)+(0,2)$}{5}
  \plG{e4}{e4 -| cA}{4}
  \plG{e12}{e12 -| cA}{12}
  \plG{e23}{e23 -| cA}{23}
  \plG{e24}{e24 -| cA}{24}
  \plG{e10}{$(e10)+(30,40)$}{10}
  \plG{e6}{$(e6)+(25,40)$}{6}
  \plG{e7}{$(e7 -| cA2)+(0,12)$}{7}
  \plG{e8}{e8 -| cA2}{8}
  \plG{e14}{e14 -| cA2}{14}
  \plG{e11}{$(e11)+(55,-18)$}{11}
  \plG{e9}{$(e9)+(-38,30)$}{9}
  \plG{e27}{$(e27 -| cB2)+(0,-20)$}{27}
  \plG{e19}{e19 -| cB2}{19}
  \plG{e28}{e28 -| cB2}{28}
  \plG{e26}{e26 -| cC}{26}
  \plG{e17}{$(e17 -| cC)+(0,6)$}{17}
  \plG{e18}{$(e18 -| cC)+(0,-8)$}{18}
  \plG{e16}{e16 -| cC}{16}
  \plG{e15}{e15 -| cC}{15}
  \plG{e13}{e13 -| cC}{13}
  \plG{e25}{e25 -| cC}{25}
  \node[font=\small\bfseries] at (0,10) {A: cabeza del rotor};
  \node[font=\small\bfseries] at (165,10) {B: cuerpo};
  \node[font=\small\bfseries] at (300,0) {C: módulo interno};
  \rotulo{$(current bounding box.south east)+(-45,-40)$}{Plano 11 -- Vista explosionada}{1:2,5}{11}
\end{tikzpicture}
\caption{Plano 11: vista explosionada en tres grupos de montaje (escala 1:2,5, isométrica).
A: cabeza del rotor, sobre el eje fijo; B: cuerpo con los patines guía propuestos
(violeta); C: módulo interno que entra por abajo. Una sola pala (8) se muestra separada
de su brazo, con la pestaña de punta (14) abajo. El orden de montaje está en la
tabla~\ref{tab:planos-montaje}.}
\label{fig:planos-explo}
\end{figure}

{\small
\begin{longtable}{@{}c>{\raggedright}p{4.0cm}c>{\raggedright}p{6.0cm}>{\raggedright\arraybackslash}p{2.9cm}@{}}
\caption{Lista de piezas del conjunto (configuración A). «Grupo» indica el grupo de
montaje de la figura~\ref{fig:planos-explo}; la masa es la del renglón de la
tabla~\ref{tab:masa} que contiene la pieza.}\label{tab:planos-piezas}\\
\toprule
\textbf{N.º} & \textbf{Pieza} & \textbf{Cant.} & \textbf{Material o componente} & \textbf{Grupo / masa}\\\midrule
\endfirsthead
\toprule
\textbf{N.º} & \textbf{Pieza} & \textbf{Cant.} & \textbf{Material o componente} & \textbf{Grupo / masa}\\\midrule
\endhead
\midrule\multicolumn{5}{r}{\footnotesize continúa}\\
\endfoot
\bottomrule
\endlastfoot
1  & Cuerda del drogue & 1 & Dyneema Ø\,\SI{1.5}{mm}, \SI{1.5}{m} & A / 18 (con 2, 20, 21)\\
2  & Copa del drogue & 1 & PETG Ø\,44 × 25, pared 1,5 & A / 18\\
3  & Tuerca de retención & 1 & M8 × 0,75, acero, traba química & A / 24 (cubo)\\
4  & Cubo giratorio & 1 & PETG o nylon, Ø\,44 × 16,5, 4 brazos & A / 24\\
5  & Rodamiento & 2 & 688ZZ ($8\times16\times5$) & A / 24\\
6  & Pasador de bisagra & 4 & Acero Ø\,2 (templado, ISO 8734, cap.~\ref{cap:estructura}) & A / 24\\
7  & Portapala & 4 & PETG impreso, paso \ang{-6} a \ang{0}, con llave & A / 30 (con 8)\\
8  & Pala & 4 & Carbono \SI{0.5}{mm} curvado 7\,\%, $156\times45$ & A / 30\\
9  & Cuerpo & 1 & PETG Ø\,88, pared \SI{2}{mm}, largo 205 & B / 141\\
10 & Sensor de rpm & 1+1 & Hall en la tapa + imán Ø\,3 en el cubo & A / 24\\
11 & Cámara cenit & 1 & Módulo de cámara, fuera del eje & A / 20 (con 18)\\
12 & Eje fijo hueco & 1 & Aluminio 6061, Ø\,8 × 5, largo 67 & A / 24\\
13 & Batería & 1 & Celda 18650 Li-ion con portapila y brida & C / 55\\
14 & Pestaña de punta & 4 & Acero inoxidable \SI{0.5}{mm} con ojal & A / 16 (traba)\\
15 & Anillo de traba y servo & 1 & PETG, 4 dientes; micro servo MG90S & C / 16\\
16 & Lastre & 1 & Disco de acero, ajustable & C / 83\\
17 & Baliza & 1 & Zumbador con pila propia (\Req{E8}) & C / 8\\
18 & Cámara nadir & 1 & Módulo de cámara, mira hacia abajo & C / 20\\
19 & Orificios de interruptores & 2 & Ø\,\SI{8}{mm} (\Req{E5}), S1 y S2 & B / ---\\
20 & Drogue & 1 & Tela de nailon ripstop, Ø\,250 & A / 18\\
21 & Destorcedor & 1 & De bolas, en la cuerda (no dibujado) & A / 18\\
22 & Separador de rodamientos & 1 & Tubo de latón o aluminio, \SI{3}{mm} & A / 24\\
23 & Tapa superior & 1 & PETG, con soporte del hall y ventana de la cámara & A / 25 (tapas)\\
24 & Travesaño de anclaje & 1 & PETG o aluminio, con argolla para la cuerda & A / 25\\
25 & Pila de placas & 3 & PCB FR-4 Ø\,76 sobre separadores M2,5 & C / 47\\
26 & Tapa inferior & 1 & PETG, apoyo axial en el cohete & C / 25\\
27 & Patín guía (\textbf{propuesto}) & 4 & PETG, $4\times3{,}5$, define el Ø\,95 & B / $\approx$\,4 (nuevo)\\
28 & Indicador de encendido & 1 & LED (\Req{E6}) visible en un hueco & B / 25 (cables)\\
29 & Resorte de torsión & 4 & \SI{0.005}{N.m} de apertura (no dibujado) & A / 24\\
30 & Paneles solares & 2 & En los huecos traseros (\Req{E3}; no se ven) & B / 8\\
\end{longtable}}

\begin{table}[H]
\centering\small
\caption{Orden de montaje propuesto. Cada paso se puede verificar antes del siguiente.}
\label{tab:planos-montaje}
\begin{tabularx}{\textwidth}{@{}clY@{}}
\toprule
\textbf{Paso} & \textbf{Grupo} & \textbf{Operación y verificación}\\\midrule
1 & C & Fijar baliza (17) y cámara nadir (18) a la tapa inferior (26); atornillar el lastre (16) y pesar el módulo.\\
2 & C & Montar servo y anillo (15) en su soporte; con el servo energizado, verificar que el giro de \ang{15} saca los cuatro dientes.\\
3 & C & Batería (13) con brida y pila de placas (25) sobre separadores; prueba eléctrica en banco.\\
4 & B+C & Introducir el módulo C por abajo en el cuerpo (9) y atornillar la tapa inferior; verificar acceso a S1 y S2 por los orificios (19).\\
5 & A & Fijar eje (12) al travesaño (24) y a la tapa superior (23); pasar la cuerda (1) por el eje y anudarla a la argolla.\\
6 & A & Rodamientos (5) y separador (22) en el cubo (4); calzar el cubo en el eje y apretar la tuerca (3). El cubo debe girar libre $\ge\SI{10}{s}$ con un impulso a mano.\\
7 & A & Portapalas (7) con palas (8), pasadores (6) y resortes (29); comprobar la llave de orientación y el paso.\\
8 & A+B & Atornillar la tapa superior al cuerpo; plegar las palas y trabar las pestañas (14) con el servo en posición de traba.\\
9 & A & Plegar el drogue (20) en la copa (2); verificar el destorcedor (21).\\
10 & Todo & Pesar: \SI{500 \pm 10}{g} (\Req{S1}); medir Ø y largo con calibre pasa/no pasa de Ø\,95,0 y \SI{250}{mm}.\\
\bottomrule
\end{tabularx}
\end{table}

La lista separa lo que viene de la base (piezas 1 a 19, iguales a las del plano~1) de lo
que este capítulo agrega para cerrar el conjunto: separador, tapas, travesaño, pila de
placas y, como propuesta, los patines guía (27), que se justifican en la
sección~\ref{sec:planos-tubo}.

%=====================================================================
\section{Plano 12: secuencia de despliegue a escala}\label{sec:planos-secuencia}

La secuencia completa tiene seis configuraciones mecánicas distintas, y en cada una la
fuerza dominante es diferente: la reacción del cohete, el tirón del drogue, el arrastre
del drogue, el aire que abre las palas, el par aerodinámico que acelera el rotor y,
finalmente, el empuje del rotor. La figura~\ref{fig:planos-secuencia} dibuja las seis a
escala 1:8 con las fuerzas a escala \SI{1}{N}\,$\widehat{=}$\,\SI{2.5}{mm}; las fuerzas
mayores de \SI{6}{N} se dibujan cortadas y se indica su valor.

Las fuerzas de las configuraciones estables salen del equilibrio vertical
(ecuación~\ref{eq:equilibrio}) con los coeficientes de la base:
\begin{align}
  \text{etapa 1:}\quad & W = D_\text{d}+D_\text{c} = \tfrac12\rho V_1^2\,(C_{D,\text{d}}A_\text{d}+C_{D,\text{c}}A_\text{c})\notag\\
     &\quad\Rightarrow\; V_1=\SI{13.35}{m/s},\; D_\text{d}=\SI{4.29}{N},\; D_\text{c}=\SI{0.62}{N},
  \label{eq:planos-eq1}\\
  \text{etapa 2:}\quad & W = T + D_\text{d}+D_\text{c}\notag\\
     &\quad\Rightarrow\; V_2=\SI{6.40}{m/s},\; T=\SI{3.78}{N},\; D_\text{d}=\SI{0.98}{N},\; D_\text{c}=\SI{0.14}{N}.
  \label{eq:planos-eq2}
\end{align}
La tensión en la cuerda es el arrastre del drogue menos su propio peso ($\approx\SI{0.12}{N}$):
\SI{4.17}{N} en la etapa 1 y \SI{0.87}{N} en la etapa 2. El arranque del rotor se estima con
un par que cae linealmente desde $Q_0=\SI{0.026}{N.m}$ (paso \ang{-2}, $\Omega=0$, ecuación~\ref{eq:Q0})
hasta cero en $\Omega_f=\SI{120}{rad/s}$:
\begin{equation}
  I_p\,\dot\Omega = Q_0\left(1-\frac{\Omega}{\Omega_f}\right)
  \;\Rightarrow\;
  \Omega(t)=\Omega_f\left(1-e^{-t/\tau}\right),\qquad
  \tau=\frac{I_p\,\Omega_f}{Q_0}=\frac{\num{5.13e-4}\times120}{0{,}026}=\SI{2.4}{s}.
  \label{eq:planos-arranque}
\end{equation}
El rotor llega al 90\,\% de su velocidad en $t_{90}=\tau\ln10\approx\SI{5.5}{s}$. Es una
estimación de orden de magnitud: el par real depende de la velocidad de descenso, que cae
mientras el rotor acelera, y del paso; el capítulo~\ref{cap:teoria} da la base del par de
arranque.

\newcommand{\plSqCuerpo}{%
  \filldraw[fijo] (-44,0) rectangle (44,205);
  \filldraw[rot] (-22,205) rectangle (22,222);
  \filldraw[fijo] (-22,225) rectangle (22,250);
  \draw[eje] (0,-15) -- (0,262);}
\newcommand{\plSqPlegadas}{%
  \filldraw[rot] (-46.5,57) rectangle (-44.5,211);
  \filldraw[rot] (44.5,57) rectangle (46.5,211);}
% pala lateral: lado (+1/-1), ángulo (0 = horizontal, -90 = colgando)
\newcommand{\plSqPala}[2]{%
  \filldraw[rot] ({#1*45.2},213) -- ++({#1*156*cos(#2)},{156*sin(#2)})
     -- ++({-#1*4*sin(#2)},{4*cos(#2)}) -- ({#1*45.2-#1*4*sin(#2)},{213+4*cos(#2)}) -- cycle;
  \fill[black] ({#1*45.2},213) circle (2.5);}
\newcommand{\plSqDrogue}{%
  \filldraw[cuerda, fill=ccuerda!15] (-125,400) .. controls (-95,465) and (95,465) .. (125,400)
     .. controls (60,388) and (-60,388) .. (-125,400);
  \foreach \x in {-125,-62,62,125}{\draw[ccuerda, line width=0.3pt] (\x,400) -- (0,330);}}
\newcommand{\plSqCuerda}{%
  \draw[cuerda] (0,250) -- (0,280); \draw[cuerda] (0,298) -- (0,330);
  \draw[black, line width=0.4pt] (-9,275) -- (9,283) (-9,293) -- (9,301);}
\newcommand{\plF}[4]{% inicio, fuerza en N (con signo), color, etiqueta a la derecha
  \draw[-{Latex[length=2.2mm]}, line width=1.1pt, #3] (#1) -- ++(0,{20*#2})
     node[pos=0.55, right, font=\scriptsize, #3, align=left] {#4};}
\newcommand{\plFl}[4]{% igual, etiqueta a la izquierda
  \draw[-{Latex[length=2.2mm]}, line width=1.1pt, #3] (#1) -- ++(0,{20*#2})
     node[pos=0.55, left, font=\scriptsize, #3, align=right] {#4};}
\newcommand{\plMarco}[4]{% título, estado, velocidad, texto inferior
  \draw[cgris!50, line width=0.3pt] (-215,-190) rectangle (215,530);
  \node[font=\small\bfseries, anchor=north west] at (-210,526) {#1};
  \node[font=\scriptsize, anchor=north east, cgris] at (210,490) {#3};
  \node[font=\scriptsize, anchor=north west, align=left, text width=50mm] at (-208,-32) {#4};
  \node[font=\scriptsize\ttfamily, anchor=south, fill=cgris!10, inner sep=2pt] at (0,-185) {#2};}

\begin{figure}[p]
\centering
\begin{tikzpicture}[x=0.125mm, y=0.125mm]
% ---------------------------------------------------------------- (a)
\begin{scope}[shift={(0,0)}]
  \plMarco{(a) En el cohete}{LAUNCH\_PAD / ASCENT}{$V$: 0 $\uparrow$ ascenso}
    {Apoyo en la tapa inferior: \SI{74}{N} a 15\,G (\Req{S4}).\\ Traba cerrada (rojo).
     Tubo Ø\,98 supuesto (sec.~\ref{sec:planos-tubo}).}
  \fill[cgris!15] (-49,-25) rectangle (-55,300) (49,-25) rectangle (55,300);
  \draw[cgris, line width=0.5pt] (-49,-25) -- (-49,300) (49,-25) -- (49,300);
  \filldraw[fill=cgris!35, draw=cgris] (-49,-22) rectangle (49,-2);
  \fill[pattern=north east lines, pattern color=cgris] (-49,-22) rectangle (49,-2);
  \plSqCuerpo\plSqPlegadas
  \fill[cfreno] (-40,58) rectangle (40,62);
  \plF{-150,150}{3.5}{cgris}{$a$}
\end{scope}
% ---------------------------------------------------------------- (b)
\begin{scope}[shift={(445,0)}]
  \plMarco{(b) Expulsión y apertura}{APOGEE $\rightarrow$ DESCENT}{$V$: 15--\SI{35}{m/s} $\downarrow$}
    {Tirón de apertura en la argolla del travesaño, no en los rodamientos. Traba cerrada.}
  \plSqCuerpo\plSqPlegadas
  \plSqCuerda
  \draw[cuerda, fill=ccuerda!15] (-60,400) .. controls (-40,465) and (40,465) .. (60,400)
     .. controls (30,372) and (-30,372) .. (-60,400);
  \foreach \x in {-60,60}{\draw[ccuerda, line width=0.3pt] (\x,400) -- (0,330);}
  \draw[-{Latex[length=2.2mm]}, line width=1.1pt, cfreno] (0,300) -- ++(0,28);
  \node[font=\scriptsize, cfreno, anchor=north west, align=left, inner sep=1pt] at (66,452)
     {tirón (fuera\\ de escala):\\ \SI{8}{N} a \SI{15}{m/s}\\ \SI{23}{N} a \SI{25}{m/s}\\
      \SI{44}{N} a \SI{35}{m/s}\\ diseño \SI{147}{N}\\ (30\,G)};
\end{scope}
% ---------------------------------------------------------------- (c)
\begin{scope}[shift={(890,0)}]
  \plMarco{(c) Etapa 1: drogue}{DESCENT}{$V_1\approx\SI{13.4}{m/s}$ $\downarrow$}
    {Equilibrio \eqref{eq:planos-eq1}: $T_c=D_\text{d}-m_\text{d}g$.\\ Duración $\approx\SI{10.5}{s}$ (apogeo \SI{700}{m}).}
  \plSqCuerpo\plSqPlegadas
  \plSqCuerda\plSqDrogue
  \plF{0,250}{4.17}{ccuerda!70!black}{$T_c=\SI{4.17}{N}$}
  \plFl{-60,107}{-4.905}{black}{$W=$\\ \SI{4.91}{N}}
  \fill[black] (0,107) circle (5);
  \plF{60,-14}{0.62}{cfijo}{$D_\text{c}=\SI{0.62}{N}$}
\end{scope}
% ---------------------------------------------------------------- (d)
\begin{scope}[shift={(0,-760)}]
  \plMarco{(d) Liberación}{PROBE\_RELEASE}{$V\approx\SI{13.4}{m/s}$ $\downarrow$}
    {El servo gira el anillo \ang{15}; aire y resorte abren las palas en 130--\SI{150}{ms};
     golpe en el tope a 41--\SI{52}{rad/s} (cap.~\ref{cap:estructura}).}
  \plSqCuerpo
  \draw[crot, densely dashed] (-46.5,57) rectangle (-44.5,211) (44.5,57) rectangle (46.5,211);
  \plSqPala{1}{-40}\plSqPala{-1}{-40}
  \draw[crot, -{Latex[length=2mm]}, line width=0.7pt] (45.2,100) arc (-90:-48:113);
  \draw[crot, -{Latex[length=2mm]}, line width=0.7pt] (-45.2,100) arc (-90:-132:113);
  \plSqCuerda\plSqDrogue
  \plF{0,250}{4.17}{ccuerda!70!black}{$T_c\approx\SI{4.2}{N}$}
\end{scope}
% ---------------------------------------------------------------- (e)
\begin{scope}[shift={(445,-760)}]
  \plMarco{(e) Arranque del rotor}{PAYLOAD\_RELEASE}{$V$: 13,4 $\to$ \SI{6.4}{m/s} $\downarrow$}
    {$Q_0=\SI{0.026}{N.m}$, $\tau\approx\SI{2.4}{s}$ \eqref{eq:planos-arranque}.\\
     $\Omega$: 0 $\to$ \SI{120}{rad/s} en $\approx\SI{5.5}{s}$.\\ El sensor hall confirma el giro.}
  \plSqCuerpo
  \plSqPala{1}{8}\plSqPala{-1}{8}
  \plSqCuerda\plSqDrogue
  \draw[crot!80!black, -{Latex[length=2mm]}, line width=0.8pt] (82,262) arc (70:110:240 and 30);
  \plF{130,237}{1.9}{crot!80!black}{$T$: 0$\to$\\ \SI{3.8}{N}}
\end{scope}
% ---------------------------------------------------------------- (f)
\begin{scope}[shift={(890,-760)}]
  \plMarco{(f) Etapa 2: autorrotación}{PAYLOAD\_RELEASE $\rightarrow$ LANDED}{$V_2\approx\SI{6.4}{m/s}$ $\downarrow$}
    {Equilibrio \eqref{eq:planos-eq2}; $\Omega\approx\SI{120}{rad/s}$, $\beta_0\approx\ang{4.4}$.\\
     Duración $\approx\SI{87}{s}$ (apogeo \SI{700}{m}).}
  \plSqCuerpo
  \plSqPala{1}{4.4}\plSqPala{-1}{4.4}
  \plSqCuerda\plSqDrogue
  \draw[crot!80!black, -{Latex[length=2mm]}, line width=0.8pt] (82,262) arc (70:110:240 and 30);
  \plF{0,302}{0.87}{ccuerda!70!black}{$T_c=\SI{0.87}{N}$}
  \plF{130,237}{3.78}{crot!80!black}{$T=$\\ \SI{3.78}{N}}
  \plFl{-60,107}{-4.905}{black}{$W=$\\ \SI{4.91}{N}}
  \fill[black] (0,107) circle (5);
  \plF{60,-14}{0.5}{cfijo}{$D_\text{c}=\SI{0.14}{N}$}
  \draw[-{Latex[length=2mm]}, line width=0.9pt, crot!80!black] (-150,196) -- (-205,196)
     node[below right, font=\scriptsize] {$F_c\approx\SI{13}{N}$};
\end{scope}
  \rotulo{$(current bounding box.south east)+(0,-110)$}{Plano 12 -- Secuencia de despliegue}{1:8}{12}
\end{tikzpicture}
\caption{Plano 12: secuencia de despliegue a escala 1:8, cuerda del drogue cortada (largo
real \SI{1.5}{m}). Fuerzas a escala \SI{1}{N}\,$\widehat{=}$\,\SI{2.5}{mm}, salvo el tirón
de (b), que está fuera de escala, y $D_\text{c}$ de (f), agrandada para que se vea. Verde:
tensión de la cuerda; naranja: rotor; azul: arrastre del cuerpo; negro: peso en el centro
de masa ($z\approx\SI{107}{mm}$). Debajo de cada cuadro, el estado del software de vuelo
(anexo). Valores estimados, no medidos.}
\label{fig:planos-secuencia}
\end{figure}

\begin{table}[H]
\centering\small
\caption{Secuencia de despliegue: disparador, duración y carga dominante. Duraciones para un
apogeo de \SI{700}{m} (tabla de tiempos del capítulo~\ref{cap:dim}).}
\label{tab:planos-secuencia}
\begin{tabularx}{\textwidth}{@{}clYYc@{}}
\toprule
\textbf{Cuadro} & \textbf{Estado} & \textbf{Disparador} & \textbf{Carga dominante} & \textbf{Duración}\\\midrule
a & \texttt{ASCENT} & Despegue (aceleración) & Inercia sobre la tapa inferior: \SI{74}{N} a 15\,G & $\approx$\,10--\SI{15}{s}\\
b & \texttt{APOGEE} & Carga de expulsión del cohete & Tirón del drogue en la argolla: 8--\SI{44}{N} & $<\SI{1}{s}$\\
c & \texttt{DESCENT} & Drogue abierto & Tensión de cuerda \SI{4.2}{N} & $\approx\SI{10.5}{s}$\\
d & \texttt{PROBE\_RELEASE} & Altura $\le80$\,\% del apogeo & Golpe de las palas en el tope (cap.~\ref{cap:estructura}) & 0,13--\SI{0.15}{s}\\
e & \texttt{PAYLOAD\_RELEASE} & Hall: $\Omega$ supera un umbral & Par aerodinámico $Q_0=\SI{0.026}{N.m}$ & $\approx\SI{5.5}{s}$\\
f & \texttt{LANDED} al tocar & Rotor en régimen & Empuje \SI{3.78}{N}; centrífuga \SI{13}{N} por pala & $\approx\SI{87}{s}$\\
\bottomrule
\end{tabularx}
\end{table}

\begin{nota}[Transición: qué falta saber]
El cuadro (e) es el menos conocido. Mientras el rotor acelera, el CanSat sigue bajando a
casi \SI{13}{m/s} y recorre del orden de 40 a \SI{60}{m} hasta que $V$ se acerca a
\SI{6.4}{m/s}. Ese tramo cuenta en el promedio de \Req{C4}; el modelo lineal
\eqref{eq:planos-arranque} no lo resuelve con precisión. El ensayo de caída con el rotor
debe medir $\Omega(t)$ con el sensor hall desde la liberación.
\end{nota}

%=====================================================================
\section{Plano 13: recorrido de cargas}\label{sec:planos-cargas}

Cada carga tiene un camino corto y conocido hasta el punto donde se reacciona, y ninguna
pasa por una pieza que no fue pensada para ella: el tirón del drogue no toca los
rodamientos, y el empuje del rotor no pasa por la traba. Esta sección lo verifica con un
modelo cuasiestático de barra a lo largo del eje y lo dibuja en la
figura~\ref{fig:planos-cargas}.

\subsection{Modelo}
El CanSat se trata como una barra en $z$ (de \SIrange{0}{250}{mm}) con las masas $m_i$ de
la tabla~\ref{tab:masa} ubicadas en su altura $z_i$ (el tubo de \SI{141}{g} se reparte en
41 tramos) y con fuerzas externas $F_k$ aplicadas en $z_k$: el empuje del rotor en el cubo
($z=\SI{213}{mm}$), la cuerda en la argolla ($z=\SI{178}{mm}$), el arrastre del cuerpo y
la reacción del cohete o del suelo en el fondo ($z=0$). Con un factor de carga
$n=(g+a)/g$ común a todas las masas, el esfuerzo axial en un corte $z$ (tracción
positiva) es
\begin{equation}
  N(z)=\sum_{z_k>z}F_k\;-\;n\,g\sum_{z_i>z}m_i ,
  \qquad n=\frac{\sum_k F_k}{m\,g}.
  \label{eq:planos-axial}
\end{equation}
Las fuerzas de interfaz de la tabla~\ref{tab:planos-interfaces-cargas} salen de aplicar
el mismo equilibrio a cada pieza. Los casos son los de la
tabla~\ref{tab:planos-secuencia} más los de diseño de \Req{S4}, \Req{S5} y del
aterrizaje del capítulo~\ref{cap:estructura}:
ascenso a 15\,G; apertura del drogue a \SI{35}{m/s} (\SI{44.2}{N}, $n=10{,}1$) y con
\SI{147}{N} ($n=30{,}8$, 30\,G); etapas 1 y 2 estables ($n=1$); y aterrizaje con
espuma a \SI{200}{g}. El modelo no incluye cargas laterales (péndulo, ráfagas, cuerda
desviada); esas están en el capítulo~\ref{cap:estructura}.

\newcommand{\plCgBase}[1]{% 1: plegado (0) o desplegado (1)
  \fill[ffijo] (-44,0) rectangle (-42,205) (42,0) rectangle (44,205);
  \draw[cfijo, line width=0.4pt] (-44,0) rectangle (-42,205) (42,0) rectangle (44,205);
  \filldraw[fijo, line width=0.4pt] (-42,0) rectangle (42,3);
  \filldraw[fijo, line width=0.4pt] (-42,199) rectangle (-4,205) (4,199) rectangle (42,205);
  \filldraw[fijo, line width=0.4pt] (-20,176) rectangle (20,180);
  \filldraw[fijo, line width=0.4pt] (-4,180) rectangle (-2.5,249) (2.5,180) rectangle (4,249);
  \filldraw[rot, line width=0.4pt] (-22,205.5) rectangle (-8,222) (8,205.5) rectangle (22,222);
  \draw[black, fill=white, line width=0.3pt] (-8,207) rectangle (-4,212) (4,207) rectangle (8,212)
        (-8,215) rectangle (-4,220) (4,215) rectangle (8,220);
  \filldraw[fill=cgris!40, draw=cgris, line width=0.3pt] (-7,222) rectangle (-4,225) (4,222) rectangle (7,225);
  \draw[cfijo, line width=0.4pt] (-22,250) -- (-22,225) -- (22,225) -- (22,250);
  \filldraw[fijo, line width=0.4pt] (-9,72) rectangle (9,137);
  \fill[cgris!40] (-36,20) rectangle (36,30);
  \foreach \yp in {146,158,170}{\draw[ccuerda!60!black, line width=0.8pt] (-38,\yp) -- (38,\yp);}
  \filldraw[act, line width=0.4pt] (-13,44) rectangle (13,66);
  \ifnum#1=0
    \filldraw[rot, line width=0.4pt] (-46,57) rectangle (-44.6,211) (44.6,57) rectangle (46,211);
  \else
    \filldraw[rot, line width=0.4pt] (-22,212) -- (-95,218) -- (-95,221) -- (-22,215) -- cycle;
    \filldraw[rot, line width=0.4pt] (22,212) -- (95,218) -- (95,221) -- (22,215) -- cycle;
    \draw[black, line width=0.4pt] (-99,214) -- (-91,225) (99,214) -- (91,225);
  \fi
  \fill[ccuerda] (0,178) circle (1.8);
  \draw[eje] (0,-12) -- (0,262);}
\tikzset{plcam/.style={line width=2.6pt, draw=#1, opacity=0.55, line cap=round, line join=round,
          -{Triangle[length=2.6mm, width=3.2mm]}},
         plval/.style={font=\scriptsize, fill=white, inner sep=1.5pt, align=left},
         plvalr/.style={font=\scriptsize, fill=white, inner sep=1.5pt, align=right}}

\begin{figure}[p]
\centering
\begin{tikzpicture}[x=0.25mm, y=0.25mm]
% ------------------------------------------------------------------ A: ascenso
\begin{scope}
  \node[font=\small\bfseries, anchor=west] at (-150,300) {(a) Ascenso, 15\,G (\Req{S4})};
  \plCgBase{0}
  \fill[cgris!35] (-49,-14) rectangle (49,-2);
  \fill[pattern=north east lines, pattern color=cgris] (-49,-14) rectangle (49,-2);
  \draw[plcam=cfreno] (15,214) -- (6,206) -- (3.2,195) -- (3.2,183) -- (25,178) -- (43,170) -- (43,8) -- (43,-10);
  \draw[plcam=cfreno] (-30,25) -- (-30,4) -- (-30,-10);
  \draw[plcam=cfreno] (-5,100) -- (-38,100) -- (-43,90) -- (-43,-10);
  \draw[ref] (12,210) -- (65,232) node[plval, anchor=west] {I4: \SI{7.9}{N}\\ hombro inferior};
  \draw[ref] (20,178) -- (65,185) node[plval, anchor=west] {I2: \SI{11.5}{N}};
  \draw[ref] (43,120) -- (65,120) node[plval, anchor=west] {pared: $-\SI{44}{N}$\\ (a media altura)};
  \draw[ref] (-20,100) -- (-65,120) node[plvalr, anchor=east] {I8 batería:\\ \SI{8.1}{N}};
  \draw[ref] (-30,15) -- (-65,40) node[plvalr, anchor=east] {I9 lastre:\\ \SI{12.2}{N}};
  \draw[ref] (43,-8) -- (65,-20) node[plval, anchor=west] {I11 apoyo: \SI{73.6}{N}};
  \draw[flecha, cgris] (-110,160) -- ++(0,50) node[midway, left, font=\scriptsize] {$a$};
\end{scope}
% ------------------------------------------------------------------ B: apertura del drogue
\begin{scope}[shift={(330,0)}]
  \node[font=\small\bfseries, anchor=west] at (-150,300) {(b) Apertura del drogue};
  \plCgBase{0}
  \draw[cuerda] (0,178) -- (0,285);
  \draw[plcam=ccuerda!80!black] (0,182) -- (0,282);
  \draw[plcam=cfreno] (-30,24) -- (-30,4) -- (-43,4) -- (-43,165) -- (-25,178) -- (-4,178);
  \draw[plcam=cfreno] (14,214) -- (6,206) -- (3.2,195) -- (3.2,183);
  \draw[ref] (0,262) -- (-40,270) node[plvalr, anchor=east] {I1 cuerda: \SI{44.2}{N}\\ (30\,G: \SI{147}{N})};
  \draw[ref] (-22,178) -- (-65,200) node[plvalr, anchor=east] {I2: \SI{37.7}{N}\\ (30\,G: \SI{127}{N})};
  \draw[ref] (-43,110) -- (-65,110) node[plvalr, anchor=east] {pared en tracción:\\ \SI{34}{N} (\SI{116}{N})};
  \draw[ref] (-30,15) -- (-65,25) node[plvalr, anchor=east] {I9: \SI{8.2}{N} (\SI{25}{N})};
  \draw[ref] (12,210) -- (65,232) node[plval, anchor=west] {I4: \SI{5.4}{N} (\SI{16}{N})\\ cubo apoya \textbf{abajo}};
  \draw[ref] (4,130) -- (65,140) node[plval, anchor=west] {no pasa por\\ los rodamientos};
\end{scope}
% ------------------------------------------------------------------ C: etapa 2
\begin{scope}[shift={(0,-400)}]
  \node[font=\small\bfseries, anchor=west] at (-150,300) {(c) Etapa 2 estable};
  \plCgBase{1}
  \draw[cuerda] (0,178) -- (0,285);
  \draw[plcam=crot] (-85,222) -- (-85,235) -- (-30,226) -- (-14,222) -- (-6,214);
  \draw[plcam=crot] (85,222) -- (85,235) -- (30,226) -- (14,222) -- (6,214);
  \draw[plcam=crot] (5.5,224) -- (3.2,215) -- (3.2,186);
  \draw[plcam=ccuerda!80!black] (0,280) -- (0,186);
  \draw[plcam=cfreno] (8,178) -- (40,172) -- (43,160) -- (43,105) -- (12,105);
  \draw[ref] (-85,232) -- (-110,250) node[plvalr, anchor=east] {$T=\SI{3.78}{N}$};
  \draw[ref] (70,219) -- (95,250) node[plval, anchor=west] {I3: $F_c\approx\SI{13}{N}$/pala\\ (se anula en el cubo)};
  \draw[ref] (5.5,223.5) -- (65,195) node[plval, anchor=west] {I4: \SI{3.25}{N}\\ contra la tuerca};
  \draw[ref] (0,262) -- (-40,280) node[plvalr, anchor=east] {I1: \SI{0.87}{N}};
  \draw[ref] (-15,178) -- (-65,170) node[plvalr, anchor=east] {I2: \SI{4.0}{N}};
  \draw[ref] (43,135) -- (65,135) node[plval, anchor=west] {pared: $\le\SI{3.6}{N}$\\ en tracción};
  \draw[ref] (-9,100) -- (-65,95) node[plvalr, anchor=east] {pesos propios:\\ batería \SI{0.54}{N}};
\end{scope}
% ------------------------------------------------------------------ D: aterrizaje
\begin{scope}[shift={(330,-400)}]
  \node[font=\small\bfseries, anchor=west] at (-150,300) {(d) Aterrizaje, 200\,g};
  \plCgBase{1}
  \fill[ccuerda!25] (-44,-25) rectangle (44,0);
  \node[font=\tiny, ccuerda!40!black] at (0,-12) {espuma 25};
  \draw[cgris, line width=0.8pt] (-80,-25) -- (80,-25);
  \fill[pattern=north east lines, pattern color=cgris] (-80,-33) rectangle (80,-25);
  \draw[plcam=cfreno] (43,-35) -- (43,2) -- (43,170) -- (25,178) -- (3.2,183) -- (3.2,195) -- (6,206) -- (12,212);
  \draw[plcam=cfreno] (-30,-35) -- (-30,4) -- (-30,19);
  \draw[plcam=cfreno] (-43,-35) -- (-43,95) -- (-12,100);
  \draw[ref] (43,-30) -- (65,-40) node[plval, anchor=west] {I11 suelo: \SI{957}{N}};
  \draw[ref] (43,100) -- (65,100) node[plval, anchor=west] {pared: $-\SI{567}{N}$\\ (media altura)};
  \draw[ref] (20,178) -- (65,180) node[plval, anchor=west] {I2: \SI{129}{N}};
  \draw[ref] (10,210) -- (65,232) node[plval, anchor=west] {I4: \SI{106}{N}\\ hombro inferior};
  \draw[ref] (-20,100) -- (-65,130) node[plvalr, anchor=east] {I8 batería:\\ \SI{108}{N}};
  \draw[ref] (-30,10) -- (-65,40) node[plvalr, anchor=east] {I9 lastre:\\ \SI{163}{N}};
\end{scope}
  \rotulo{$(current bounding box.south east)+(0,-60)$}{Plano 13 -- Recorrido de cargas}{1:4}{13}
\end{tikzpicture}
\caption{Plano 13: recorrido de cargas en cuatro casos (corte simplificado a escala 1:4).
Las franjas indican el camino de la carga y su sentido de transmisión; los números son
fuerzas en cada interfaz (tabla~\ref{tab:planos-interfaces-cargas}). En (b), entre
paréntesis, el caso de diseño de 30\,G (\Req{S5}). Palas desplegadas recortadas en (c) y
(d). Valores estimados con \eqref{eq:planos-axial}.}
\label{fig:planos-cargas}
\end{figure}

\begin{figure}[H]
\centering
\begin{tikzpicture}
\begin{axis}[width=0.47\textwidth, height=7.2cm, xmin=-12, xmax=38, ymin=0, ymax=250,
  xlabel={$N(z)$ [N] (tracción $+$)}, ylabel={$z$ [mm]},
  legend style={font=\scriptsize, at={(0.98,0.02)}, anchor=south east}, title={\small (a) Vuelo}]
  \addplot[cfreno, thick] table[x=dro35, y=z] {analysis/data/planos_axial.dat};
  \addlegendentry{apertura, \SI{35}{m/s}}
  \addplot[ccuerda!70!black, thick, dashed] table[x=eta1, y=z] {analysis/data/planos_axial.dat};
  \addlegendentry{etapa 1}
  \addplot[crot, very thick] table[x=eta2, y=z] {analysis/data/planos_axial.dat};
  \addlegendentry{etapa 2}
  \draw[black!50] (axis cs:0,0) -- (axis cs:0,250);
  \draw[cgris, densely dotted] (axis cs:-12,178) -- (axis cs:38,178) node[pos=0.02, above right, font=\tiny] {argolla};
  \draw[cgris, densely dotted] (axis cs:-12,213) -- (axis cs:38,213) node[pos=0.02, above right, font=\tiny] {cubo};
\end{axis}
\end{tikzpicture}\hfill
\begin{tikzpicture}
\begin{axis}[width=0.47\textwidth, height=7.2cm, xmin=-110, xmax=130, ymin=0, ymax=250,
  xlabel={$N(z)$ [N] (tracción $+$)}, ylabel={$z$ [mm]},
  legend style={font=\scriptsize, at={(0.98,0.98)}, anchor=north east}, title={\small (b) Diseño}]
  \addplot[cgris, thick] table[x=asc15, y=z] {analysis/data/planos_axial.dat};
  \addlegendentry{ascenso 15\,G}
  \addplot[cfreno, thick] table[x=dro30G, y=z] {analysis/data/planos_axial.dat};
  \addlegendentry{apertura 30\,G}
  \addplot[cfijo, very thick, dashed] table[x expr=\thisrow{ater200}/10, y=z] {analysis/data/planos_axial.dat};
  \addlegendentry{aterrizaje 200\,g ($\div10$)}
  \draw[black!50] (axis cs:0,0) -- (axis cs:0,250);
  \draw[cgris, densely dotted] (axis cs:-110,178) -- (axis cs:130,178);
\end{axis}
\end{tikzpicture}
\caption{Esfuerzo axial a lo largo del eje del CanSat, ecuación~\eqref{eq:planos-axial}
(datos en \texttt{planos\_axial.dat}). El salto en $z=\SI{178}{mm}$ es la cuerda; el de
$z=\SI{213}{mm}$, el empuje del rotor. Arriba de la argolla la pieza trabaja en
compresión al abrir el drogue: el cubo y la copa apoyan hacia abajo.}
\label{fig:planos-axial}
\end{figure}

\begin{table}[H]
\centering\small
\caption{Fuerzas en las interfaces principales [N], salida de
\texttt{planos\_cargas.py}. Las de \SI{200}{g} coinciden con las del
capítulo~\ref{cap:estructura}.}
\label{tab:planos-interfaces-cargas}
\begin{tabularx}{\textwidth}{@{}lYrrrrr@{}}
\toprule
\textbf{Id.} & \textbf{Interfaz} & \textbf{15\,G} & \textbf{35\,m/s} & \textbf{30\,G} & \textbf{Etapa 2} & \textbf{200\,g}\\\midrule
I1  & Cuerda -- argolla del travesaño        & 0    & 44,2 & 147  & 0,87 & 0\\
I2  & Travesaño -- pared                      & 11,5 & 37,7 & 127  & 4,0  & 129\\
I3  & Pasador de bisagra (axial de la pala)   & 1,1  & 0,7  & 2,2  & 13,0 & 14,5\\
I4  & Cubo -- rodamientos -- eje (axial)      & 7,9  & 5,4  & 16,3 & 3,25 & 106\\
I8  & Batería -- brida                       & 8,1  & 5,5  & 16,6 & 0,54 & 108\\
I9  & Lastre -- tapa inferior                 & 12,2 & 8,2  & 25,0 & 0,81 & 163\\
I10 & Cuerpo, máximo $|N(z)|$                 & 73,6 & 34,0 & 116  & 3,6  & 957\\
I11 & Apoyo inferior (cohete o suelo)         & 73,6 & ---  & ---  & ---  & 957\\
\bottomrule
\end{tabularx}
\end{table}

\begin{resultado}[Lo que fija el recorrido de cargas]
\begin{enumerate}[leftmargin=*]
  \item \textbf{El aterrizaje dimensiona las fijaciones internas}, no la apertura ni el
        ascenso: la tapa inferior y la pared reciben 13 veces la carga de 15\,G. Las
        uniones de batería, lastre y rotor deben verificarse a \SI{110}{N},
        \SI{165}{N} y \SI{106}{N}.
  \item \textbf{El cubo necesita un hombro inferior} en el eje: al abrir el drogue y al
        aterrizar el cubo empuja hacia abajo (\SIrange{5}{106}{N}); la tuerca solo trabaja
        en la etapa 2 (\SI{3.25}{N}).
  \item \textbf{La pared trabaja en tracción} entre la tapa inferior y el travesaño durante
        la apertura del drogue: las uniones tapa--pared deben ser tornillos con insertos,
        no un simple encastre a presión.
  \item La centrífuga (\SI{13}{N} por pala) es la mayor carga permanente del rotor, pero se
        anula dentro del cubo y no llega al cuerpo; solo un desbalance la transmite
        (\SIrange{0.17}{1.7}{N} para 0,1 a \SI{1}{g}, capítulo~\ref{cap:dim}).
\end{enumerate}
\end{resultado}

El rotor también transmite un par al cuerpo a través de los rodamientos. Con un
coeficiente de fricción de catálogo para rodamientos rígidos de bolas,
$\mu\approx0{,}0015$, el par por carga axial es
$M=\tfrac12\mu T d=\tfrac12\times0{,}0015\times3{,}78\times0{,}008\approx\SI{2.3e-5}{N.m}$,
al que se suma el roce de la grasa y de los escudos, del mismo orden o mayor. Es un par
pequeño, pero no tiene contrapartida: el cuerpo tenderá a girar lentamente en el sentido
del rotor hasta que su propio arrastre de giro lo equilibre. El destorcedor de la cuerda
evita que ese giro se transmita al drogue; las aletas antigiro opcionales del plano~2 lo
limitarían si las cámaras lo necesitan.

%=====================================================================
\section{Plano 14: diagrama N$^2$ de interfaces}\label{sec:planos-n2}

El diagrama N$^2$ muestra que el conjunto tiene 31 interfaces (14 mecánicas, 7 eléctricas y
10 de datos) y que dos subsistemas las concentran: la estructura entrega 8 y la aviónica
intercambia 12 (6 de entrada y 6 de salida). En la matriz de la figura~\ref{fig:planos-n2} cada subsistema ocupa una casilla de
la diagonal; la casilla de la fila $i$ y la columna $j$ contiene lo que $i$ entrega a $j$
(convención usual de la ingeniería de sistemas: salidas por las filas, entradas por las
columnas). Las casillas vacías indican que no hay interfaz directa. La
tabla~\ref{tab:planos-n2} da la magnitud de cada interfaz y cómo se verifica; es la tabla
de referencias de este plano.

Los subsistemas son diez. El cohete aparece como subsistema externo porque impone
tres interfaces críticas: el apoyo, la expulsión y la holgura del tubo. Las piezas se agrupan
así: drogue y cuerda (1, 20, 21); rotor (3 a 8, 22, 29); traba (14, 15); estructura (2, 9,
12, 16, 19, 23, 24, 26, 27); potencia (13, interruptores S1 y S2, indicador 28, paneles
30); aviónica (microcontrolador, barómetro, IMU, GPS, sensor hall 10, memoria); cámaras (11,
18); radio y estación terrena; baliza (17).

\newcommand{\plNdiag}[2]{%
  \filldraw[fill=cfijo!18, draw=cfijo!70!black, line width=0.5pt]
    ({#1-0.5},{#1-0.5}) rectangle ({#1+0.5},{#1+0.5});
  \node[font=\scriptsize\bfseries, align=center, text width=1.35cm] at (#1,#1) {#2};}
% casilla: fila, columna, tipo (M/E/D), texto
\newcommand{\plNc}[4]{%
  \ifx M#3\def\plNcol{cfijo}\fi
  \ifx E#3\def\plNcol{cfreno}\fi
  \ifx D#3\def\plNcol{ccont}\fi
  \fill[\plNcol!12] ({#2-0.5},{#1-0.5}) rectangle ({#2+0.5},{#1+0.5});
  \node[font=\scriptsize\bfseries, text=\plNcol!85!black, align=center] at (#2,#1) {#4};}

\begin{figure}[H]
\centering
\begin{tikzpicture}[x=1.48cm, y=-1.48cm]
  % casillas con interfaz
  \plNc{1}{2}{M}{M2}  \plNc{1}{3}{M}{M3}  \plNc{1}{5}{M}{M1}  \plNc{1}{7}{D}{D6}
  \plNc{2}{3}{M}{M6}  \plNc{2}{5}{M}{M4}
  \plNc{3}{5}{M}{M7}  \plNc{3}{7}{D}{D1}  \plNc{3}{8}{D}{D7}  \plNc{3}{9}{D}{D10}
  \plNc{4}{3}{M}{M9}
  \plNc{5}{2}{M}{M5}  \plNc{5}{3}{M}{M8}  \plNc{5}{4}{M}{M10} \plNc{5}{6}{M}{M11}
  \plNc{5}{8}{M}{M13} \plNc{5}{10}{M}{M14}
  \plNc{6}{4}{E}{E2}  \plNc{6}{7}{E}{E1}  \plNc{6}{8}{E}{E3}  \plNc{6}{9}{E}{E4}
  \plNc{7}{6}{E}{E5}  \plNc{7}{9}{D}{D2}
  \plNc{9}{7}{D}{D3}  \plNc{10}{9}{D}{D8}
  % casillas dobles (eléctrica + datos) y la interfaz de presión estática
  \fill[cfreno!12] (3.5,6.5) rectangle (4.5,7.5);
  \fill[ccont!12] (3.5,7) rectangle (4.5,7.5);
  \node[font=\scriptsize\bfseries, text=cfreno!85!black] at (4,6.78) {E7};
  \node[font=\scriptsize\bfseries, text=ccont!85!black] at (4,7.25) {D4};
  \fill[cfreno!12] (7.5,6.5) rectangle (8.5,7.5);
  \fill[ccont!12] (7.5,7) rectangle (8.5,7.5);
  \node[font=\scriptsize\bfseries, text=cfreno!85!black] at (8,6.78) {E6};
  \node[font=\scriptsize\bfseries, text=ccont!85!black] at (8,7.25) {D5};
  \fill[ccont!12] (6.5,4.5) rectangle (7.5,5.5);
  \node[font=\scriptsize\bfseries, text=ccont!85!black] at (7,4.78) {D9};
  \node[font=\scriptsize\bfseries, text=cfijo!85!black] at (7,5.25) {M12};
  % celda prohibida: potencia -> baliza
  \fill[pattern=north east lines, pattern color=cgris!70] (9.5,5.5) rectangle (10.5,6.5);
  \node[font=\tiny, fill=white, inner sep=1pt, align=center] at (10,6) {\Req{E8}\\ prohibida};
  % rejilla
  \foreach \i in {0,...,10}{
    \draw[cgris!60, line width=0.3pt] (0.5,{\i+0.5}) -- (10.5,{\i+0.5});
    \draw[cgris!60, line width=0.3pt] ({\i+0.5},0.5) -- ({\i+0.5},10.5);}
  % diagonal
  \plNdiag{1}{Cohete\\ \mdseries(externo)}
  \plNdiag{2}{Drogue y cuerda}
  \plNdiag{3}{Rotor}
  \plNdiag{4}{Traba}
  \plNdiag{5}{Estruc\-tura}
  \plNdiag{6}{Potencia}
  \plNdiag{7}{Aviónica}
  \plNdiag{8}{Cámaras}
  \plNdiag{9}{Radio y tierra}
  \plNdiag{10}{Baliza}
  % cadena crítica de liberación (lazo cerrado)
  \draw[cfreno, line width=2.4pt, opacity=0.35, rounded corners=4pt]
    (4.4,6.6) -- (4.4,3.6) -- (2.6,3.6) -- (2.6,2.6) -- (7.4,2.6) -- (7.4,6.6) -- cycle;
  \foreach \a/\b in {{(4.4,5.1)}/{(4.4,4.9)}, {(3.6,3.6)}/{(3.4,3.6)}, {(5.4,2.6)}/{(5.6,2.6)}, {(7.4,4.9)}/{(7.4,5.1)}, {(5.6,6.6)}/{(5.4,6.6)}}{
    \draw[cfreno, -{Latex[length=2.6mm]}, line width=0.8pt] \a -- \b;}
  % flechas de convención
  \draw[flecha, cgris] (0.5,0.2) -- (2.2,0.2) node[right, font=\scriptsize] {entradas por columna};
  \draw[flecha, cgris] (0.2,0.5) -- (0.2,2.2);
  \node[font=\scriptsize, cgris, rotate=-90, anchor=west] at (0.2,2.3) {salidas por fila};
  % leyenda
  \begin{scope}[shift={(0.5,11.0)}]
    \fill[cfijo!12] (0,-0.15) rectangle (0.3,0.15); \node[right, font=\scriptsize] at (0.35,0) {M: mecánica (14)};
    \fill[cfreno!12] (2.7,-0.15) rectangle (3.0,0.15); \node[right, font=\scriptsize] at (3.05,0) {E: eléctrica (7)};
    \fill[ccont!12] (5.1,-0.15) rectangle (5.4,0.15); \node[right, font=\scriptsize] at (5.45,0) {D: datos o señal (10)};
    \draw[cfreno, line width=2.4pt, opacity=0.35] (0,0.6) -- (0.3,0.6);
    \node[right, font=\scriptsize] at (0.35,0.6) {cadena de liberación del rotor: D4/E7 $\to$ M9 $\to$ D1};
  \end{scope}
  \rotulo{10.5,12.0}{Plano 14 -- Diagrama N$^2$}{sin escala}{14}
\end{tikzpicture}
\caption{Plano 14: diagrama N$^2$ de interfaces. Fila $i$, columna $j$: lo que el
subsistema $i$ entrega al $j$. La casilla rayada (potencia $\to$ baliza) es una interfaz
que el anexo prohíbe (\Req{E8}). La banda roja sigue la cadena que libera el rotor: la
aviónica ordena al servo (D4, E7), la traba suelta las palas (M9) y el sensor hall del
rotor confirma el giro a la aviónica (D1). Códigos en la tabla~\ref{tab:planos-n2}.}
\label{fig:planos-n2}
\end{figure}

{\small
\begin{longtable}{@{}l>{\raggedright}p{1.9cm}>{\raggedright}p{6.2cm}>{\raggedright\arraybackslash}p{5.3cm}@{}}
\caption{Interfaces del diagrama N$^2$ (tabla de referencias del plano 14). Las fuerzas
vienen de la tabla~\ref{tab:planos-interfaces-cargas} y son estimaciones; las
verificaciones son ensayos propuestos, no realizados.}\label{tab:planos-n2}\\
\toprule
\textbf{Id.} & \textbf{De $\to$ a} & \textbf{Interfaz} & \textbf{Magnitud y verificación}\\\midrule
\endfirsthead
\toprule
\textbf{Id.} & \textbf{De $\to$ a} & \textbf{Interfaz} & \textbf{Magnitud y verificación}\\\midrule
\endhead
\midrule\multicolumn{4}{r}{\footnotesize continúa}\\
\endfoot
\bottomrule
\endlastfoot
\multicolumn{4}{@{}l}{\textit{Mecánicas}}\\
M1  & Cohete $\to$ estructura & Apoyo axial en la tapa inferior (26) y guía radial por los patines (27) & \SI{73.6}{N} a 15\,G (I11); fricción de extracción $\le\SI{0.21}{N}$. Calibre anillo Ø\,95,0\\
M2  & Cohete $\to$ drogue & Expulsión: el viento relativo saca el drogue de la copa sin orden del software (\Req{C2}) & Ensayo de extracción en un tramo de tubo\\
M3  & Cohete $\to$ rotor & Tubo frente a las palas plegadas & Holgura en los bordes $\ge\SI{0.2}{mm}$ con patines (sec.~\ref{sec:planos-tubo})\\
M4  & Drogue $\to$ estructura & Cuerda por el eje hueco hasta la argolla del travesaño (24) & \SI{44}{N} a \SI{35}{m/s}; \SI{147}{N} de diseño (\Req{S5}). Tracción a $1{,}5\times\SI{147}{N}$\\
M5  & Estructura $\to$ drogue & Copa (2) y eje hueco (12) alojan el drogue y guían la cuerda & Copa Ø\,41 interior × 25. Plegado repetible\\
M6  & Drogue $\to$ rotor & La cuerda pasa por el eje fijo y no toca el cubo & Cuerda Ø\,1,5 en eje Ø\,5: \SI{1.75}{mm} radial. Cubo libre con cuerda tensa\\
M7  & Rotor $\to$ estructura & Empuje por rodamientos (5) y tuerca (3); par de fricción & \SI{3.25}{N} en la tuerca (I4); $\approx\SI{2e-5}{N.m}$ más grasa. Giro libre $\ge\SI{10}{s}$\\
M8  & Estructura $\to$ rotor & Hombro inferior del eje: apoyo del cubo en ascenso, apertura y aterrizaje & \SI{16}{N} a 30\,G; \SI{106}{N} a \SI{200}{g} (I4)\\
M9  & Traba $\to$ rotor & Dientes del anillo (15) retienen las pestañas (14) & \SI{1.1}{N} lateral por pala a 15\,G (I3). Liberación repetida en banco\\
M10 & Estructura $\to$ traba & Soporte del servo y del anillo; ranuras de las pestañas en la pared & Giro de \ang{15} sin trabarse con el cuerpo cargado\\
M11 & Estructura $\to$ potencia & Brida de la batería (13); orificios de S1 y S2 (19) & \SI{8.1}{N} a 15\,G; \SI{108}{N} a \SI{200}{g} (I8). Acceso a S1 y S2 (\Req{E5})\\
M12 & Estructura $\to$ aviónica & Pila de placas (25) sobre separadores M2,5 & \SI{6.9}{N} a 15\,G; \SI{92}{N} a \SI{200}{g} (\SI{47}{g})\\
M13 & Estructura $\to$ cámaras & Soportes y ventanas de las cámaras (11, 18) & \SI{1.5}{N} a 15\,G; \SI{20}{N} a \SI{200}{g} cada una\\
M14 & Estructura $\to$ baliza & Fijación en la tapa inferior; acceso a su interruptor & \SI{16}{N} a \SI{200}{g}. Interruptor accesible (\Req{E9})\\
\multicolumn{4}{@{}l}{\textit{Eléctricas}}\\
E1  & Potencia $\to$ aviónica & \SIlist{3.3;5}{V} regulados & Caída de tensión con el servo activo\\
E2  & Potencia $\to$ traba & \SI{5}{V} al servo MG90S & Picos de arranque: medir con osciloscopio\\
E3  & Potencia $\to$ cámaras & Alimentación de las dos cámaras & Autonomía $\ge\SI{2}{h}$ en el cohete (\Req{E7})\\
E4  & Potencia $\to$ radio & Alimentación de la radio & Picos de transmisión\\
E5  & Aviónica $\to$ potencia & Medición de tensión y corriente & Telemetría (\Req{SN3}, \Req{SN8})\\
E6  & Aviónica $\to$ cámaras & Habilitación de la alimentación & Encendido desde el software\\
E7  & Aviónica $\to$ traba & Señal PWM del servo & \SI{50}{Hz}; posición de traba y de liberación\\
\multicolumn{4}{@{}l}{\textit{Datos y señales}}\\
D1  & Rotor $\to$ aviónica & Imán del cubo frente al hall (10): rpm & 1 pulso por vuelta, $\approx\SI{19}{Hz}$ en régimen; confirma \texttt{PAYLOAD\_RELEASE}\\
D2  & Aviónica $\to$ tierra & Telemetría & \SI{1}{Hz} (\Req{C5})\\
D3  & Tierra $\to$ aviónica & Comandos del anexo & Prueba de enlace en tierra\\
D4  & Aviónica $\to$ traba & Orden de liberación, estado \texttt{PROBE\_RELEASE} & Al \SI{80}{\%} de $h_\text{max}$ (\Req{C4})\\
D5  & Aviónica $\to$ cámaras & Inicio y fin de grabación & Registro en la memoria de cada cámara\\
D6  & Cohete $\to$ aviónica & Aceleración y presión del ascenso: detección de lanzamiento y apogeo & Barómetro a \SI{50}{Hz}; simulación con datos de vuelo\\
D7  & Rotor $\to$ cámaras & Las palas cruzan el campo de la cámara cenit & Paso de palas \SI{76.5}{Hz}; alias de \SI{1.5}{Hz} a 25\,fps y \SI{13.5}{Hz} a 30\,fps\\
D8  & Baliza $\to$ tierra & Señal audible para la recuperación & Alcance audible en campo\\
D9  & Estructura $\to$ aviónica & Puerto estático del barómetro (orificio en la pared) & Lejos de la estela del rotor y de los orificios de S1 y S2\\
D10 & Rotor $\to$ radio & Palas de carbono (conductoras) cerca de las antenas & Medir el enlace con las palas plegadas y abiertas\\
\end{longtable}}

Tres observaciones salen directamente de la matriz.

\begin{enumerate}[leftmargin=*]
  \item \textbf{La función principal atraviesa cuatro subsistemas.} La liberación del rotor
        es la cadena aviónica $\to$ traba $\to$ rotor $\to$ aviónica (D4/E7, M9, D1). Es un
        lazo cerrado: si el hall no ve giro después de la orden, el software puede repetir
        la orden de liberación. Cualquier ensayo de liberación debe probar la cadena
        completa y no cada eslabón por separado.
  \item \textbf{La estructura es la fila más llena.} Entrega siete interfaces
        mecánicas y una de datos; casi todas las fuerzas de la tabla~\ref{tab:planos-interfaces-cargas}
        pasan por ella. Cambiar la pared o las tapas obliga a revisar M1, M4, M8 y M10 a
        M14.
  \item \textbf{Hay interfaces no evidentes.} D7 (el rotor en el campo de la cámara
        cenit) produce un efecto estroboscópico: con $\Omega\approx\SI{120}{rad/s}$ y
        cuatro palas el paso es de \SI{76.5}{Hz}, y a 25 cuadros por segundo las palas
        parecerán girar a \SI{1.5}{Hz}. D10 (carbono cerca de las antenas) puede atenuar la
        radio y el GPS; la magnitud depende de dónde vayan las antenas, que aún no está
        definido. Ninguna de las dos aparece en los planos del capítulo~\ref{cap:planos}.
\end{enumerate}

\begin{nota}[Interfaz prohibida]
La baliza no puede tomar energía de la batería principal (\Req{E8}) y su interruptor debe
ser accesible desde afuera (\Req{E9}). En la matriz eso aparece como una casilla rayada en
potencia $\to$ baliza y como la interfaz mecánica M14. Ningún cable debe unir la baliza al
resto, ni siquiera una masa común.
\end{nota}

%=====================================================================
\section{Plano 15: envolvente en el tubo del cohete}\label{sec:planos-tubo}

Con las palas tal como están en el capítulo~\ref{cap:planos}, sus bordes definen el
Ø\,\SI{95}{mm} y son la primera pieza que toca el tubo del cohete: la holgura en los bordes
después del contacto es exactamente cero. Con cuatro patines guía de Ø\,95 y las palas
reducidas a Ø\,\SI{94}{mm} en los bordes, el contacto lo toman siempre los patines y a las
palas les quedan al menos \SI{0.22}{mm}. Esta sección dibuja el conjunto dentro del tubo
(figuras~\ref{fig:planos-tubo-long} y~\ref{fig:planos-tubo-aa}) y calcula las holguras.

\subsection{Datos y suposiciones}

El anexo fija el diámetro del CanSat (Ø\,\SI{95}{mm}, \Req{S2}), su largo
(\SIrange{200}{250}{mm}, \Req{S3}), que va en la parte superior del cohete (\Req{C1}),
que sale con la carga de expulsión (\Req{C2}) y que debe deslizar libremente (\Req{S7}).
\textbf{No da el diámetro interior del tubo, el largo de la bahía ni cómo es el
apoyo.} Los planos de esta sección usan las suposiciones siguientes, que hay que
confirmar con la organización:

\begin{itemize}[leftmargin=*]
  \item tubo de Ø\,\SI{98}{mm} interior nominal con tolerancia de \SI{\pm0.25}{mm} y pared
        de \SI{2}{mm}, de fibra de vidrio, cartón fenólico o aluminio;
  \item el CanSat apoya su tapa inferior sobre un protector de gases de \SI{14}{mm} y una
        mampara o pistón; arriba queda una holgura axial de \SI{15}{mm} hasta el hombro de
        la ojiva (bahía de al menos \SI{265}{mm});
  \item la copa del drogue mira hacia la ojiva, de modo que el drogue sale primero;
  \item temperatura en la rampa entre \SIlist{5;45}{\degreeCelsius} (montaje a
        \SI{20}{\degreeCelsius}).
\end{itemize}

\newcommand{\plC}[3]{\draw[ref] #1 -- (#2); \node[marca] at (#2) {#3};}

\begin{figure}[p]
\centering
\begin{tikzpicture}[x=0.5mm, y=0.5mm]
% ================================================================ corte longitudinal 1:2
\begin{scope}
  % tubo del cohete (31)
  \foreach \s in {-1,1}{
    \filldraw[fill=cgris!15, draw=cgris!80!black, line width=0.4pt] ({\s*49},-42) rectangle ({\s*51},300);
    \fill[pattern=north east lines, pattern color=cgris] ({\s*49},-42) rectangle ({\s*51},300);}
  \draw[cgris!80!black, line width=0.4pt, decorate, decoration={zigzag, segment length=6, amplitude=1.5}]
    (-56,300) -- (56,300) (-56,-42) -- (56,-42);
  % hombro de la ojiva (32), mampara inferior (34) y protector de gases (33)
  \filldraw[fill=cgris!35, draw=cgris!80!black, line width=0.4pt] (-49,265) rectangle (49,288);
  \filldraw[fill=cgris!35, draw=cgris!80!black, line width=0.4pt] (-49,-30) rectangle (49,-14);
  \fill[pattern=north east lines, pattern color=cgris!80!black] (-49,-30) rectangle (49,-14);
  \filldraw[fill=crot!8, draw=cgris, line width=0.3pt] (-49,-14) rectangle (49,0);
  \fill[pattern=dots, pattern color=cgris] (-49,-14) rectangle (49,0);
  % gases de expulsión
  \foreach \x in {-30,0,30}{\draw[cfreno, -{Latex[length=2mm]}, line width=0.7pt] (\x,-56) -- (\x,-33);}
  \node[font=\scriptsize, cfreno, anchor=north] at (0,-56) {gases de la carga de expulsión (\Req{C2})};
  % CanSat
  \filldraw[fijo] (-44,0) rectangle (44,205);
  \filldraw[fill=cfijo!8, draw=cfijo!60, line width=0.3pt] (-42,2) rectangle (42,203);
  \filldraw[fill=cfijo!30, draw=cfijo!70!black, line width=0.3pt] (-9.3,72) rectangle (9.3,137);
  \foreach \zz in {150,158,166}{\draw[ccuerda!70!black, line width=0.9pt] (-38,\zz) -- (38,\zz);}
  \filldraw[fill=cgris!50, draw=cgris!80!black, line width=0.3pt] (-38,18) rectangle (38,32);
  \filldraw[rot] (-22,205) rectangle (22,222);
  \filldraw[fill=cgris!40, draw=cgris!80!black, line width=0.3pt] (-7,222) rectangle (7,225);
  \filldraw[fijo] (-22,225) rectangle (22,250);
  \filldraw[fill=ccuerda!25, draw=ccuerda!70!black, line width=0.3pt] (-20.5,227) rectangle (20.5,248.5);
  \draw[eje] (0,-48) -- (0,295);
  % palas plegadas en el plano de corte (8) y patines fuera del plano (27)
  \foreach \s in {-1,1}{
    \fill[crot] ({\s*44.4},57) rectangle ({\s*45.2},211);
    \draw[ccont, densely dashed, line width=0.4pt] ({\s*44},0) rectangle ({\s*47.5},205);}
  % envolvente Ø95
  \draw[cgris, densely dashed, line width=0.3pt] (-47.5,-8) -- (-47.5,258) (47.5,-8) -- (47.5,258);
  % sección A-A
  \foreach \s in {-1,1}{
    \draw[black, line width=1.1pt] ({\s*56},130) -- ({\s*66},130);
    \draw[black, -{Latex[length=2mm]}, line width=0.6pt] ({\s*63},130) -- ({\s*63},120);
    \node[font=\small\bfseries, anchor=south] at ({\s*61},131) {A};}
  % cotas
  \draw[cota] (-49,318) -- (49,318) node[midway, fill=white, font=\scriptsize] {Ø\,98 int. (supuesto)};
  \draw[cgris, line width=0.3pt] (-49,300) -- (-49,322) (49,300) -- (49,322);
  \draw[cota] (-47.5,-76) -- (47.5,-76) node[midway, fill=white, font=\scriptsize] {Ø\,95 (\Req{S2})};
  \draw[cgris, line width=0.3pt] (-47.5,-8) -- (-47.5,-80) (47.5,-8) -- (47.5,-80);
  \draw[cota] (-80,0) -- (-80,250) node[midway, fill=white, font=\scriptsize, rotate=90] {250 (\Req{S3})};
  \draw[cgris, line width=0.3pt] (-22,250) -- (-84,250) (-49,0) -- (-84,0);
  \draw[cota] (-80,250) -- (-80,265) node[midway, left, font=\scriptsize] {15};
  \draw[cgris, line width=0.3pt] (-49,265) -- (-84,265);
  % llamadas
  \plC{(49,280)}{80,280}{31}
  \plC{(30,276)}{80,262}{32}
  \plC{(30,-7)}{80,-7}{33}
  \plC{(30,-22)}{80,-22}{34}
  \plC{(15,240)}{80,240}{2}
  \plC{(45.2,190)}{80,200}{8}
  \plC{(47.5,100)}{80,100}{27}
  \plC{(42,60)}{80,60}{9}
  \plC{(-35,25)}{-100,25}{16}
  \plC{(-6,110)}{-100,110}{13}
  \plC{(-30,158)}{-100,158}{25}
  \plC{(-15,213)}{-100,213}{4}
\end{scope}
\rotulo{195,-85}{Plano 15 -- CanSat en el tubo (hoja 1)}{1:2}{15}
\end{tikzpicture}
\caption{Plano 15, hoja 1: corte longitudinal del CanSat plegado dentro del tubo del cohete
(escala 1:2). Las palas (8) están cortadas por
su plano medio; los patines (27) están a \ang{45} de las palas, fuera del plano de corte
(línea de trazos violeta). El tubo (31), la ojiva (32), el protector de gases (33) y la
mampara (34) son \textbf{supuestos}: el anexo no los define. Referencias en la
tabla~\ref{tab:planos-ref15}.}
\label{fig:planos-tubo-long}
\end{figure}

\begin{figure}[p]
\centering
\begin{tikzpicture}[x=1mm, y=1mm]
  % tubo (31)
  \filldraw[fill=cgris!15, draw=cgris!80!black, line width=0.4pt, even odd rule] (0,0) circle (51) (0,0) circle (49);
  \fill[pattern=north east lines, pattern color=cgris, even odd rule] (0,0) circle (51) (0,0) circle (49);
  % envolvente Ø95
  \draw[cgris, densely dashed, line width=0.3pt] (0,0) circle (47.5);
  % cuerpo (9)
  \filldraw[fill=ffijo, draw=cfijo, line width=0.5pt, even odd rule] (0,0) circle (44) (0,0) circle (42);
  \fill[cfijo!5] (0,0) circle (42);
  % batería (13) y separadores de la pila de placas (25)
  \filldraw[fill=cfijo!30, draw=cfijo!70!black, line width=0.4pt] (0,0) circle (9.3);
  \foreach \a in {45,135,225,315}{\filldraw[fill=cgris!50, draw=cgris!80!black, line width=0.3pt] (\a:34) circle (1.6);}
  \draw[cfijo!60, line width=0.3pt] (-9.3,4) -- (-30,4) (9.3,4) -- (30,4) (-9.3,-4) -- (-30,-4) (9.3,-4) -- (30,-4);
  % patines (27)
  \foreach \a in {45,135,225,315}{
    \begin{scope}[rotate=\a-90]
      \filldraw[fill=ccont!30, draw=ccont!80!black, line width=0.4pt] (-2,43.9) rectangle (2,47.5);
    \end{scope}}
  % palas plegadas (8), bordes en Ø94
  \foreach \a in {0,90,180,270}{
    \begin{scope}[rotate=\a-90]
      \draw[crot, line width=1.4pt] (0,-37.48) ++(74.05:81.9) arc (74.05:105.95:81.9);
    \end{scope}}
  % detalle X
  \draw[black, line width=0.4pt] (37:46.5) circle (7);
  \node[font=\small\bfseries, anchor=south west] at ($(37:46.5)+(4.5,5)$) {X};
  % cotas
  \draw[cota] (-49,-60) -- (49,-60) node[midway, fill=white, font=\scriptsize] {Ø\,98 int. (supuesto)};
  \draw[cgris, line width=0.3pt] (-49,0) -- (-49,-63) (49,0) -- (49,-63);
  \draw[cota] (-47.5,-67) -- (47.5,-67) node[pos=0.5, fill=white, font=\scriptsize] {Ø\,95 patines};
  \draw[cgris, line width=0.3pt] ({47.5*cos(225)},{47.5*sin(225)}) -- (-47.5,-70)
                                 ({47.5*cos(315)},{47.5*sin(315)}) -- (47.5,-70);
  \draw[cota] (-62,{-47*sin(28.3)}) -- (-62,{47*sin(28.3)});
  \draw[cgris, line width=0.3pt] ({-47*cos(28.3)},{47*sin(28.3)}) -- (-65,{47*sin(28.3)})
                                 ({-47*cos(28.3)},{-47*sin(28.3)}) -- (-65,{-47*sin(28.3)});
  \node[font=\scriptsize, fill=white, rotate=90] at (-62,0) {cuerda 45};
  \draw[cota] (0,0) -- (45:30);
  \node[font=\scriptsize, fill=white, inner sep=1pt] at (45:20) {\ang{45}};
  \draw[cota] (0:56) arc (0:28.3:56);
  \node[font=\scriptsize, anchor=west] at (14:57) {\ang{28.3}};
  \draw[cgris, line width=0.3pt] (0:44) -- (0:58) (28.3:47) -- (28.3:58);
  \draw[cgris, line width=0.3pt] (0,0) -- (0:44);
  % llamadas
  \plC{(100:50)}{-30,62}{31}
  \plC{(125:43)}{-58,60}{9}
  \plC{(135:46)}{-72,40}{27}
  \plC{(-6:45.2)}{66,-18}{8}
  \plC{(-5,5)}{-25,25}{13}
  \plC{(225:34)}{-60,-45}{25}
  \node[font=\scriptsize, anchor=north west, align=left] at (54,-36)
    {Ø\,94: bordes de pala\\ Ø\,88: cuerpo\\ pared \SI{2}{mm} (\Req{S6})};
% ---------------------------------------------------------------- detalle X 3:1
\begin{scope}[shift={(0,-112)}, x=3mm, y=3mm]
  \begin{scope}[shift={(0,-46.5)}, rotate=53]
    \clip (24:41) -- (24:52.5) arc (24:50:52.5) -- (50:41) arc (50:24:41) -- cycle;
    \filldraw[fill=cgris!15, draw=cgris!80!black, line width=0.4pt] (20:49) arc (20:54:49) -- (54:51) arc (54:20:51) -- cycle;
    \fill[pattern=north east lines, pattern color=cgris] (20:49) arc (20:54:49) -- (54:51) arc (54:20:51) -- cycle;
    \filldraw[fill=ffijo, draw=cfijo, line width=0.5pt] (20:42) arc (20:54:42) -- (54:44) arc (54:20:44) -- cycle;
    \begin{scope}[rotate=-45]
      \filldraw[fill=ccont!30, draw=ccont!80!black, line width=0.5pt] (-2,43.9) rectangle (2,47.5);
    \end{scope}
    \begin{scope}[rotate=-90]
      \draw[crot, line width=1.5pt] (0,-37.48) ++(74.05:81.9) arc (74.05:105.95:81.9);
    \end{scope}
    \draw[cgris, densely dashed, line width=0.3pt] (20:47.5) arc (20:54:47.5);
    % cotas radiales
    \draw[cota] (45:47.5) -- (45:49);
    \draw[cota] (28.3:47.0) -- (28.3:49);
    \coordinate (cP) at (45:48.25); \coordinate (cB) at (28.3:48);
    \coordinate (pP) at (45:45.7); \coordinate (pB) at (27:46.4);
    \coordinate (pT) at (40:50); \coordinate (pC) at (40:43);
    \coordinate (aA) at (28.3:42.5); \coordinate (aB) at (45:42.5);
  \end{scope}
  \draw[cgris, line width=0.3pt] (aA) -- ++(0,0);
  \node[font=\scriptsize, anchor=south west, fill=white, inner sep=1pt] at ($(cP)+(1.2,1.5)$) {1,5};
  \draw[ref] (cP) -- ($(cP)+(1.2,1.5)$);
  \node[font=\scriptsize, anchor=south east, fill=white, inner sep=1pt] at ($(cB)+(-1.5,2.0)$) {2,0};
  \draw[ref] (cB) -- ($(cB)+(-1.5,2.0)$);
  \node[marca] at ($(pP)+(-5.5,-1)$) {27}; \draw[ref] (pP) -- ($(pP)+(-5.5,-1)$);
  \node[marca] at ($(pB)+(4,-2.5)$) {8};  \draw[ref] (pB) -- ($(pB)+(4,-2.5)$);
  \node[marca] at ($(pT)+(4,4.5)$) {31}; \draw[ref] (pT) -- ($(pT)+(4,4.5)$);
  \node[marca] at ($(pC)+(-3,-5.5)$) {9};  \draw[ref] (pC) -- ($(pC)+(-3,-5.5)$);
  \node[font=\small\bfseries, anchor=south] at (0,9) {Detalle X (sección A-A), 3:1};
  \node[font=\scriptsize, align=left, anchor=north west, text width=75mm] at (-27,-12)
    {Holguras radiales nominales con el cuerpo centrado: \SI{1.5}{mm} en los patines y
     \SI{2.0}{mm} en los bordes de las palas. El borde de pala está a \ang{16.7} del
     patín.};
\end{scope}
  \rotulo{80,-170}{Plano 15 -- Sección A-A (hoja 2)}{1:1 y 3:1}{15}
\end{tikzpicture}
\caption{Plano 15, hoja 2: sección A-A a $z=\SI{130}{mm}$ (escala 1:1, mirando hacia abajo) y
detalle X (escala 3:1).
Cuatro patines de PETG (27) de $4\times\SI{3.5}{mm}$ a \ang{45} de las palas fijan el
Ø\,95; las palas (8) quedan con los bordes en Ø\,94 y el centro a \SI{0.42}{mm} del
cuerpo. Se ven la batería (13) y los separadores de la pila de placas (25). Abajo, el
detalle X a escala 3:1 con las holguras radiales nominales.}
\label{fig:planos-tubo-aa}
\end{figure}

\begin{table}[H]
\centering\small
\caption{Referencias del plano 15. «Origen» indica si el dato viene del anexo, del diseño o
de una suposición de este capítulo.}
\label{tab:planos-ref15}
\begin{tabularx}{\textwidth}{@{}clYl@{}}
\toprule
\textbf{N.º} & \textbf{Elemento} & \textbf{Dato} & \textbf{Origen}\\\midrule
2  & Copa del drogue & Ø\,44 × 25, hacia la ojiva & Diseño\\
4  & Cubo & Ø\,44 × 16,5 & Diseño\\
8  & Palas plegadas & Bordes en Ø\,94 (antes Ø\,95) & Diseño modificado\\
9  & Cuerpo & Ø\,88, pared \SI{2}{mm} & Diseño (\Req{S6})\\
13 & Batería & 18650, Ø\,18,6 × 65 & Diseño\\
16 & Lastre & Disco de acero & Diseño\\
25 & Pila de placas & Ø\,76; separadores M2,5 (posición dibujada en Ø\,68) & Diseño\\
27 & Patines guía & 4 × (4 × 3,5 × 205), Ø\,95 & Propuesto aquí\\
31 & Tubo del cohete & Ø\,98 interior \SI{\pm0.25}{mm}, pared 2 & \textbf{Supuesto}\\
32 & Hombro de la ojiva & Holgura axial de \SI{15}{mm} & \textbf{Supuesto}\\
33 & Protector de gases & Algodón o Nomex, \SI{14}{mm} & \textbf{Supuesto}\\
34 & Mampara o pistón & Apoyo de la tapa inferior (I11) & \textbf{Supuesto}\\
\bottomrule
\end{tabularx}
\end{table}

\subsection{Holgura diametral}

La holgura diametral en el peor caso es la diferencia entre el diámetro interior más chico
del tubo y el diámetro exterior más grande del CanSat, cada uno a la temperatura del
momento:
\begin{equation}
  c_\text{peor} = D_t\,(1+\alpha_t\,\Delta T) - \delta_t
     \;-\; \bigl[D\,(1+\alpha_p\,\Delta T) + \delta_p + o\bigr],
  \label{eq:planos-holgura}
\end{equation}
con $D=\SI{95}{mm}$, $\delta_t=\SI{0.25}{mm}$ (tolerancia del tubo),
$\delta_p=\SI{0.20}{mm}$ (tolerancia de impresión de los patines), $o=\SI{0.30}{mm}$
(ovalización del cuerpo impreso), $\alpha_p=\SI{6.8e-5}{K^{-1}}$ (PETG) y
$\alpha_t$ entre \SIlist{1.5e-5;2.5e-5}{K^{-1}} según el material del tubo. Las tolerancias son
valores típicos de impresión FDM y de tubos comerciales, no mediciones; la ecuación se
puede rehacer con las del equipo.

Con $D_t=\SI{98}{mm}$ la holgura nominal es de \SI{3.0}{mm} y la de peor caso de
\SI{2.1}{mm} en caliente (tubo de fibra de vidrio, $\Delta T=+\SI{25}{K}$); en frío sube
a \SI{2.3}{mm}, porque el PETG se contrae más que el tubo. La temperatura cambia la
holgura en menos de \SI{0.15}{mm}: lo que manda son las tolerancias. Para conservar
\SI{1}{mm} en el peor caso hace falta $D_t\ge\SI{96.9}{mm}$
(figura~\ref{fig:planos-holgura}a). El material del tubo casi no influye: las curvas de
aluminio y cartón se superponen con la de fibra de vidrio.

\subsection{Juego lateral y contacto de las palas}

El cuerpo puede moverse lateralmente dentro del tubo hasta que el primer punto de su
contorno toca la pared. El script \texttt{planos\_seccion.py} recorre todas las
direcciones $\varphi$ de ese desplazamiento con el contorno real de la sección A-A
(arcos de pala de radio \SI{81.9}{mm}, patines y cuerpo) y, para cada una, calcula el
desplazamiento hasta el contacto y la holgura que queda en los bordes de las palas
(figura~\ref{fig:planos-holgura}b):

\begin{itemize}[leftmargin=*]
  \item \textbf{Diseño base (sin patines):} el juego lateral es de \SIrange{1.5}{1.7}{mm} y
        el contacto siempre es un borde de pala. Ese borde es una lámina de carbono de
        \SI{0.5}{mm} sujeta solo por la bisagra y por la pestaña de punta; al deslizar
        puede engancharse en una unión del tubo, rasparse o empujar la pestaña.
  \item \textbf{Con patines y bordes en Ø\,94:} el juego lateral es de
        \SIrange{1.5}{2.0}{mm} y el contacto siempre es un patín. La holgura mínima en los
        bordes de pala es de \SI{0.22}{mm} (tubo de Ø\,98) y \SI{0.31}{mm} (Ø\,97), en la
        dirección $\varphi=0$, cuando dos patines tocan a la vez. Es poca, pero se debe a
        la geometría y no a las tolerancias: el patín más cercano está a solo
        \ang{16.7} del borde de la pala.
  \item Para llevar los bordes de Ø\,95 a Ø\,94 basta con acercar la pala
        \SI{0.56}{mm} al eje; el centro del arco queda a \SI{0.42}{mm} del cuerpo
        (antes \SI{0.98}{mm}). No se puede bajar más sin que la pala toque la pared.
\end{itemize}

\begin{figure}[H]
\centering
\begin{tikzpicture}
\begin{axis}[width=0.5\textwidth, height=6.2cm, xmin=95, xmax=100, ymin=-1, ymax=5.2,
  xlabel={Ø interior del tubo $D_t$ [mm]}, ylabel={Holgura diametral [mm]},
  legend style={font=\scriptsize, at={(0.02,0.98)}, anchor=north west},
  title={\small (a) Holgura frente a $D_t$}, grid=major, grid style={cgris!20}]
  \addplot[cfijo, thick] table[x=Dt, y=nom] {analysis/data/planos_holgura.dat};
  \addlegendentry{nominal, \SI{20}{\degreeCelsius}}
  \addplot[cfreno, thick] table[x=Dt, y=fv_cal] {analysis/data/planos_holgura.dat};
  \addlegendentry{peor caso, \SI{45}{\degreeCelsius}}
  \addplot[ccont, thick, dashed] table[x=Dt, y=fv_frio] {analysis/data/planos_holgura.dat};
  \addlegendentry{peor caso, \SI{5}{\degreeCelsius}}
  \addplot[crot, densely dotted, thick] table[x=Dt, y=al_cal] {analysis/data/planos_holgura.dat};
  \addlegendentry{peor, aluminio, \SI{45}{\degreeCelsius}}
  \draw[black!60, dashed] (axis cs:95,1) -- (axis cs:100,1) node[pos=0.97, above left, font=\tiny] {mínimo \SI{1}{mm}};
  \draw[black!60, dashed] (axis cs:98,-1) -- (axis cs:98,5.2) node[pos=0.04, right, font=\tiny] {supuesto};
  \draw[black!60, dashed] (axis cs:96.88,-1) -- (axis cs:96.88,1);
  \node[font=\tiny, anchor=north, fill=white, inner sep=1pt] at (axis cs:96.88,-0.2) {96,9};
\end{axis}
\end{tikzpicture}\hfill
\begin{tikzpicture}
\begin{axis}[width=0.5\textwidth, height=6.2cm, xmin=0, xmax=90, ymin=-0.05, ymax=0.8,
  xtick={0,15,...,90},
  xlabel={Dirección del desplazamiento $\varphi$ [\si{\degree}]}, ylabel={Holgura en los bordes de pala [mm]},
  legend style={font=\scriptsize, at={(0.5,0.98)}, anchor=north},
  title={\small (b) Palas después del contacto}, grid=major, grid style={cgris!20}]
  \addplot[ccont, thick] table[x=phi, y=g_pat98] {analysis/data/planos_seccion.dat};
  \addlegendentry{patines, $D_t=98$}
  \addplot[ccont!60, thick, dashed] table[x=phi, y=g_pat97] {analysis/data/planos_seccion.dat};
  \addlegendentry{patines, $D_t=97$}
  \addplot[cfreno, very thick] table[x=phi, y=g_base98] {analysis/data/planos_seccion.dat};
  \addlegendentry{base, sin patines}
\end{axis}
\end{tikzpicture}
\caption{(a) Holgura diametral con la ecuación~\eqref{eq:planos-holgura} (tubo de fibra de
vidrio salvo indicación; datos en \texttt{planos\_holgura.dat}). (b) Holgura que queda en
los bordes de las palas cuando el cuerpo, desplazado en la dirección $\varphi$ (medida
desde el centro de una pala), toca el tubo (datos en \texttt{planos\_seccion.dat}). Sin
patines la holgura es cero en todas las direcciones: el contacto es la pala.}
\label{fig:planos-holgura}
\end{figure}

\subsection{Extracción}

El CanSat no puede atascarse en el tubo por efecto cajón. Si la fuerza que lo empuja actúa a
una distancia $a$ del eje, el cuerpo se inclina, toca en dos puntos opuestos separados el
largo $L$ y aparecen dos normales $N=F\,a/L$. Hay atasco si la fricción supera al empuje:
\begin{equation}
  2\,\mu\,\frac{F\,a}{L} \ge F \quad\Longleftrightarrow\quad a \ge a_\text{crit}=\frac{L}{2\mu}.
  \label{eq:planos-cajon}
\end{equation}
Con $L=\SI{250}{mm}$ y $\mu$ entre \numlist{0.3;0.5} (PETG sobre fibra de vidrio o cartón,
valor de manual), $a_\text{crit}$ queda entre \SIlist{250;417}{mm}, mucho más que el radio del
CanSat (\SI{47.5}{mm}). Además, la inclinación máxima posible dentro del tubo es de
\ang{0.69} (resolviendo $L\sin\gamma + D\cos\gamma = D_t$), de modo que el cuerpo no puede
cruzarse.

La fuerza para sacarlo es pequeña. Con el cohete inclinado \ang{5} y todo el peso
apoyado en la pared, la fricción $\mu W\sin\ang{5}$ queda entre \SIlist{0.13;0.21}{N}. El drogue
abierto a \SI{15}{m/s} tira con $\tfrac12\rho V^2 C_D A_\text{d}\approx\SI{5.4}{N}$, 25
veces más; basta con que el drogue salga de la copa y se infle para que arrastre el
CanSat fuera del tubo.

\begin{nota}[Puntos abiertos de la envolvente]
\begin{itemize}[leftmargin=*]
  \item \textbf{Diámetro del tubo:} sin él, la holgura de \SI{2.1}{mm} es solo un
        ejemplo. Si el tubo resultara de Ø\,\SI{96}{mm} o menos, el Ø\,95 del anexo deja
        menos de \SI{1}{mm} en el peor caso y hay que ajustar los patines al tubo real.
  \item \textbf{Juego en la vibración:} el juego lateral de \SIrange{1.5}{2.0}{mm} deja
        que el CanSat golpee el tubo con la vibración de 15\,G (\Req{S4}). Los patines
        reciben esos golpes en lugar de las palas, pero conviene un ensayo de vibración
        dentro de un tramo de tubo y, después, uno de extracción.
  \item \textbf{Gases de expulsión:} la tapa inferior y la ventana de la cámara nadir
        miran a la carga de expulsión. El protector (33) es responsabilidad del cohete y
        debe acordarse con la organización.
\end{itemize}
\end{nota}

%=====================================================================
\section{Síntesis de los planos}\label{sec:planos-sintesis}

La tabla~\ref{tab:planos-indice} resume los siete planos de este capítulo y lo que
cada uno fija para la fabricación y los ensayos.

\begin{table}[H]
\centering\small
\caption{Índice de los planos de conjunto avanzados.}
\label{tab:planos-indice}
\begin{tabularx}{\textwidth}{@{}clcY@{}}
\toprule
\textbf{Plano} & \textbf{Contenido} & \textbf{Escala} & \textbf{Qué fija}\\\midrule
9  & Isométricas del plegado (fig.~\ref{fig:planos-iso-plegado}) & 1:2 & Orden de apilado y acceso a interruptores\\
10 & Isométrica desplegada (fig.~\ref{fig:planos-iso-desp}) & 1:3,5 & Disco de Ø\,400 y posición del drogue\\
11 & Vista explosionada (fig.~\ref{fig:planos-explo}) & 1:2,5 & 30 piezas en 3 grupos; orden de montaje\\
12 & Secuencia de despliegue (fig.~\ref{fig:planos-secuencia}) & 1:8 & Fuerzas y estados de software en cada fase\\
13 & Recorrido de cargas (fig.~\ref{fig:planos-cargas}) & 1:4 & Hombro inferior del eje; uniones atornilladas\\
14 & Diagrama N$^2$ (fig.~\ref{fig:planos-n2}) & s/e & 31 interfaces; cadena de liberación\\
15 & CanSat en el tubo (figs.~\ref{fig:planos-tubo-long} y~\ref{fig:planos-tubo-aa}) & 1:2, 1:1, 3:1 & Patines Ø\,95; palas a Ø\,94\\
\bottomrule
\end{tabularx}
\end{table}

\begin{resultado}[Cambios al diseño que salen de estos planos]
\begin{enumerate}[leftmargin=*]
  \item Agregar cuatro \textbf{patines guía} (27) de Ø\,95 a \ang{45} de las palas y
        llevar los bordes de las palas plegadas a Ø\,94. Sin ellos, las palas son lo
        primero que toca el tubo.
  \item Dar al eje fijo un \textbf{hombro inferior} que tome el cubo en ascenso, apertura
        y aterrizaje (hasta \SI{106}{N}); la tuerca solo trabaja en la etapa 2.
  \item Unir tapas y pared con \textbf{tornillos e insertos}: la pared trabaja en tracción
        al abrir el drogue (hasta \SI{116}{N} en el caso de diseño).
  \item Ensayar la \textbf{cadena de liberación completa} (D4/E7 $\to$ M9 $\to$ D1) y
        definir dónde van las antenas, lejos de las palas de carbono.
  \item \textbf{Confirmar con la organización} el diámetro interior del tubo, el largo de
        la bahía y el protector de gases.
\end{enumerate}
\end{resultado}
